\documentclass[12pt]{article}

% Packages
\usepackage[margin=1in]{geometry}
\usepackage{amsmath, amssymb}
\usepackage{authblk}
\usepackage{hyperref}
\usepackage{natbib}
\usepackage{setspace}
\usepackage{soul}       % for \hl{} highlighting (marks in original)
\usepackage{microtype}
\usepackage{times}
\usepackage{titlesec}
\usepackage{parskip}
\usepackage{xcolor}

\hypersetup{
    colorlinks=true,
    linkcolor=blue,
    citecolor=blue,
    urlcolor=blue
}

\doublespacing

% Highlight command to mark text that was highlighted in the original
\newcommand{\highlight}[1]{\hl{#1}}

%-----------------------------------------------------------------------
\begin{document}

%-----------------------------------------------------------------------
% TITLE
%-----------------------------------------------------------------------
\begin{titlepage}
\centering
\vspace*{2cm}

{\Large\bfseries Single-cell transcriptomics to unravel the sexual dimorphism in the aorta transcriptional response to circadian disruption}

\vspace{1.5cm}

Benjamin J. Auerbach$^{1}$, Soon Y. Tang$^{2}$, Ronan Lordan$^{2}$,
Se\'{a}n T. Anderson$^{2}$, Elizabeth Hennessy$^{2}$, Ryan McConnell$^{2}$,
Mingyao Li$^{3}$, Garret A. FitzGerald$^{2}$

\vspace{1cm}

\begin{flushleft}
$^{1}$Graduate Group in Genomics and Computational Biology, University of Pennsylvania Perelman School of Medicine, Philadelphia, PA 19104, USA\\[4pt]
$^{2}$Institute for Translational Medicine and Therapeutics, University of Pennsylvania Perelman School of Medicine, Philadelphia, PA 19104, USA\\[4pt]
$^{3}$Department of Biostatistics, Epidemiology and Informatics, University of Pennsylvania Perelman School of Medicine, Philadelphia, PA 19104, USA.
\end{flushleft}

\vspace{1cm}

\begin{flushleft}
\textbf{Correspondence:}\\
Garret A. FitzGerald\\
Institute for Translational Medicine and Therapeutics\\
University of Pennsylvania Perelman School of Medicine\\
Room 10-122, Smilow Center for Translational Research, 3400 Civic Center Boulevard\\
Philadelphia, PA 19104\\
Email: \href{mailto:garret@upenn.edu}{garret@upenn.edu}\\
Tel: (215) 898-1185\\
Fax: (215) 573-9135
\end{flushleft}

\end{titlepage}

%-----------------------------------------------------------------------
% ABSTRACT
%-----------------------------------------------------------------------
\section*{Abstract}

The circadian molecular clock is a 24-hour cellular timekeeper that influences many features of cardiovascular function. Disruption of the circadian clock via misalignment with the light-dark cycle is associated with a higher incidence of cardiovascular disease and raises cardiovascular risk factors in humans. Nonetheless, the molecular basis for the deleterious effects of misalignment remains ill defined. Moreover, bulk transcriptomics from mice suggests females are more resilient to the effects of circadian misalignment in liver. Whether a sexual dimorphism exists in the response to misalignment in the vasculature remains unknown. To address these questions, we performed several single-cell RNA-sequencing experiments from whole aorta at ZT0, ZT6, ZT12, and ZT18 over the circadian cycle. Leveraging a Bayesian statistical approach, we first defined the cell-type-specific and sex-dependent circadian transcriptomes of aortic cell types at baseline. Among several sex differences, female vascular smooth muscle cells are broadly more rhythmic than their male counterparts and demonstrate a marked reduction in heat shock protein expression at ZT12. Female smooth muscle cells also demonstrate expression alterations indicative of an increased sensitivity to phenotypic switching. We next defined the cell-type-specific and sex-dependent circadian transcriptomes of aortic cell types after inducible deletion of the core circadian clock gene \textit{Bmal1}. This suggested several features of vascular function were under direct clock control, such as vascular tone, immune cell recruitment, focal adhesion, and cellular proliferation. Additionally, loss of \textit{Bmal1} in SMCs induced a more femininized expression profile and an increase in estrogen receptor expression, suggesting sex-dependent expression may be mediated, in part, through the circadian clock. Lastly, we defined the cell-type-specific and sex-dependent circadian transcriptomes of aortic cell types after acute misalignment with the light-dark cycle. We identify several expression changes indicative of increased unfolded proteins and decreased ability to activate the unfolded protein response pathway in smooth muscle cells. Macrophage chemokine receptors \textit{Cx3cr1} and \textit{Ccr2} do not adapt their timing to the new light-dark cycle, suggesting that improper timing of myeloid cell migration may be a consequence of acute misalignment. \highlight{Further inspection suggested this was a more general phenomenon, as core clock expression in macrophages entrained more slowly than the endogenous vascular cell types.} As such, acute circadian misalignment may induce a disruption of the relative timing of the vascular cell type clocks.

\medskip
\noindent\textit{MAYBE FINISH WITH A STATEMENT AROUND IDENTIFYING NEW MECHANISMS TO INVESTIGATE\ldots AND ACKNOWLEDGE TEMPO 2}

%-----------------------------------------------------------------------
% INTRODUCTION
%-----------------------------------------------------------------------
\section*{Introduction}

The circadian molecular clock is a 24-hour timekeeping mechanism found in nearly every cell in humans. The time of the clock, referred to as circadian phase, is determined by the mRNA and protein concentrations of the clock's constituent genes, referred to as clock or core clock genes. Clock genes are organized in a transcriptional-translational feedback loop that enables cells to maintain self-sustained $\sim$24-hour oscillations in the concentrations of clock gene mRNA. Clock gene proteins additionally interact with cell-type-specific regulatory factors to drive rhythmic transcription of genes referred to as clock-controlled genes (CCGs). It is largely through these CCGs that circadian clocks generate rhythmic cellular and organismal behaviors.

Many features of cardiovascular function are under direct control of the circadian clock, such as vascular tone~\cite{ref1}, thrombus formation~\cite{ref2}, endothelial nitric oxide synthase activity~\cite{ref3}, and myeloid cell recruitment and adhesion~\cite{ref4,ref5}. Crucially, disruption of the circadian clock modulates vascular disease. Disruption of the clock via light-dark (LD) misalignment by forced cycle desynchrony protocols increases cardiovascular risk factors in humans~\cite{ref6}. Epidemiological studies associate shift work with an increased incidence and prevalence of cardiovascular disease~\cite{ref7,ref8,ref9}. Genetic ablation of the clock in mice influences atherosclerosis via multiple orthogonal mechanisms, such as effects on vascular smooth muscle cell (SMC) migration~\cite{ref10}, SMC phenotypic switching and plaque stability~\cite{ref11}, lipid metabolism~\cite{ref12}, and rhythmic myeloid recruitment to the vasculature~\cite{ref4}. Moreover, existing evidence suggests clock function of the endogenous vascular cell types restrains atherosclerosis~\cite{ref13}. Nonetheless, the molecular basis of how clock disruption affects key pathways implicated in vascular disease, and in which cell types, remains incompletely understood.

In humans, biomarkers of cardiovascular risk differ between males and females. These include arterial stiffness~\cite{ref14}, lipid profiles~\cite{ref15,ref16}, blood pressure~\cite{ref17}, and immune phenotypes~\cite{ref18}. The presentation of vascular disease also differs between sexes, such as differences atherosclerotic plaque size, composition, and stability~\cite{ref19,ref20}. Existing literature suggests females exhibit more robust circadian characteristics than males and their clocks appear more resilient to disruption~\cite{ref21}. Considering this, and given the clock's influence over key vascular disease risk factors~\cite{ref22,ref23,ref24}, characterizing how the sexes differ in their molecular response to clock perturbations will be important to understand whether sex differences in vascular disease presentation are mediated by the circadian clock.

%-----------------------------------------------------------------------
% RESULTS
%-----------------------------------------------------------------------
\section*{Results}

\subsection*{Single-cell transcriptomics to unravel the cell-type-specific and sex-specific circadian transcriptome of the aorta}

To investigate these questions, we performed several single-cell RNA-sequencing (scRNA-seq) experiments from whole aorta samples harvested over circadian time courses (\highlight{Figure 1a}). First, to assess the contributions of sex and cell type on circadian expression of cells aligned with the LD cycle, whole aorta samples were collected from 12-week-old male and female C57BL/6 mice aligned with a 12:12 LD cycle over a circadian time course (\highlight{Figure 1}b). Mice were entrained to a 24-hour LD cycle for 2 weeks. Mice were then sacrificed at Zeitgeiber time (ZT) 0, 6, 12, or 18 hours, and whole aorta samples were harvested. Second, to assess sex-specific responses to genetic perturbation of the clock, \textit{Bmal1} inducible global deletion models were generated using Cre-Lox recombination; at 8 weeks of age, \textit{Bmal1} was inducibly deleted in \textit{Bmal1}$^{\text{fl/fl}}$:CAGG$^{\text{creERt2/+}}$ mice by administering tamoxifen via oral gavage (100 mg/kg of body weight) for 5 days. At 12 weeks of age, whole aorta samples were then collected from male and female \textit{Bmal1}$^{\text{fl/fl}}$:CAGG$^{\text{creERt2/+}}$ mice aligned with a 12:12 LD cycle over a circadian time course. Mice were entrained to a 24-hour LD cycle for 2 weeks. Mice were then sacrificed at ZT 0, 6, 12, or 18 hours, and whole aorta samples were harvested. Third, to assess the contribution of sex on the response to circadian misalignment with the LD cycle, 12-week-old male and female C57BL/6 mice were entrained to a 24-hour LD cycle for 2 weeks. To induce an acute misalignment with the LD cycle, the start of the LD cycle was advanced 6 hours after the entrainment period (\highlight{Figure 1}c). Mice were then sacrificed 5 days after at ZT 0, 6, 12, or 18 hours, and whole aorta samples were harvested. For all aorta samples, cells were then disassociated into single-cell suspensions that were used as input for single-cell RNA-sequencing (scRNA-seq) using the 10X Genomics Chromium platform.

In total, we generated 152,473 cells. Cells were then quality controlled for library size (minimum of 1000 unique molecular identifiers (UMI)) and mitochondrial UMI fraction (maximum of 0.2). Cells were furthered filtered for the presence of doublets using Scrublet~\cite{ref25}. This yielded 141,752 high-quality cells. To obtain a low-dimensional embedding of the cells, 1150 highly variable genes were detected (details described in Methods); the UMI count matrix was limited to these highly variable genes and used as input to scVI~\cite{ref26}. scVI was run considering library size, \textit{Bmal1} knockout status, misalignment status, batch, sex, and time as covariates to mitigate their contribution to describing variation in the low-dimensional embedding. Indeed, inspection of the results suggested variation in the low-dimensional embedding was not described by these covariates. Using the low-dimensional embedding, cells were then clustered using the Leiden algorithm~\cite{ref27} with resolution 0.05. This yielded 5 main clusters corresponding to 115,460 vascular smooth muscle cells (SMCs), 19,759 fibroblasts, 3,935 endothelial cells (ECs), 2,917 macrophages, and 631 T cells (\highlight{Figures 1d, 1f}). As the low-dimensional representation suggested substantial heterogeneity within SMCs, SMCs were furthered subclustered using the Leiden algorithm with a resolution of 0.2. This yielded 5 SMC subclusters (\highlight{Figures 1e, 1g}). The vast majority of SMCs expressed markers associated with canonical contractile SMCs. Variation within these contractile SMCs was primarily described by physical location, as SMC0 (\textit{Hoxb5+, Hoxb6+}) and SMC1 (\textit{Lgals3}+, \textit{Hand2+}) expressed markers of the aortic arch and ascending aorta~\cite{ref28,ref29,ref30}, while SMC2 (\textit{Hoxa10+, Hoxc10+, Hoxa9+, Rgs5+}) expressed markers of the descending thoracic aorta~\cite{ref28,ref29,ref30,ref31}. In addition, we detected smaller subpopulations expressing markers consistent with previously observed synthetic SMC phenotypes; SMC3 expressed markers associated with osteoblasts~\cite{ref32,ref33} (\textit{Runx2, Frzb}) and chondrocytes~\cite{ref34} (\textit{Sox9}), while SMC4 expressed markers of multipotent SEM~\cite{ref35} cells (\textit{Ly6a, Ly6c1, Vcam1}).

Given assignments of cells to clusters, we next sought to unravel the effect of sex, misalignment status, genotype, and cell type on circadian expression. To investigate this, approximate posterior distributions of gene temporal parameters were inferred using Bayesian variational inference for each gene, cell type, and condition combination (details of which can be viewed in Methods).

\subsection*{Key cell type specific functions cycle in male aorta}

Confident circadian cycling genes were identified by decomposing the expression posterior waveforms using a Fast Fourier Transform (FFT) and then computing the expected decrease in data likelihood when assuming the 24-hour component amplitude was zero (additional details of which can be viewed in Methods). This approach suggested that hundreds of genes cycle over the circadian cycle in each major cell type and SMC subtype in male mice aligned with the LD cycle (\highlight{Figures 2a and 2b}).

As a data quality check, we first assessed if the peak times of SMC cycling genes in our single-cell data aligned well with the peak times of aorta cycling genes from bulk transcriptomics, given that SMCs comprise most of the vessel wall. The expected posterior acrophases of called circadian cyclers in SMCs align well with their peak times from bulk aorta generated by \textit{Zhang et al.}~\cite{ref36} (\highlight{Figure 2c}). As anticipated, the core circadian clock genes demonstrate strong rhythms across all major cell types (\highlight{Figure 2d}). Core circadian clock genes demonstrate strong rhythms across all SMC subtypes as well (\highlight{Supplementary Figure 1}). The acrophase distribution of circadian cyclers appears bimodal, with a peak at dawn and dusk (\highlight{Supplementary Figure 2}). Altogether, these analyses provided confidence in both our data quality, as well as our approach to calling cycling genes and parameter estimation.

We next assessed aspects of cycling expression relevant to SMC phenotypic switching in atherosclerosis. Work from \textit{Pan et al.}~\cite{ref35} and others report~\cite{ref37,ref38,ref39} that SMCs dedifferentiate into SEM cells. \highlight{SEM cells, which are multipotent, can thereafter differentiate into macrophage-like cells and fibroblast-like cells}~\cite{ref35}. Confident cyclers in SMCs that monotonically increase from SMC to SEM appear to peak more frequently at dusk (ZT12) (\highlight{Figure 4a}), and are enriched for dusk acrophases (\highlight{Figure 4b}). In contrast, confident SMC cyclers that monotonically decrease from SMC to SEM appear to more frequently peak at dawn (ZT0) (\highlight{Figure 4c}), and are enriched for dawn acrophases (\highlight{Figure 4d}). As the functional role of SEM cells in atherosclerosis continues to be illuminated, it remains unknown whether they possess functional circadian clocks. Inspection of core clock gene expression in SEM cells suggests they have functional circadian clocks, with similar amplitudes and mesors to other SMC subtypes (\highlight{Supplementary Figure 1}).

Next, we identified pathways enriched for circadian cycling genes within each cell type. To do this, we first restricted genes to those with evidence of circadian expression. To quantify this, we compute what we refer to as the circadian Bayes factor, which is defined as the fold change increase in Bayesian evidence of the original waveform relative to the evidence of the waveform with the 24-hour component subtracted. Genes were restricted to those with circadian Bayes factors of at least 10 (i.e.\ the original waveform explains the data at least 10 times better than the waveform lacking the 24-hour component). Using mouse gene sets from Kyoto Encyclopedia of Genes and Genomes (KEGG) 2019, we next restricted gene sets to those with gene expression peaking at similar times. To do this, we computed gene set acrophase posterior distributions by averaging the acrophase posterior distributions of cycling genes within the gene sets. Genes sets were then filtered for 70 percent average acrophase posterior intervals smaller than 14 hours. Thereafter, we identified overrepresented gene sets over the circadian cycle using Enrichr~\cite{ref40}.

This analysis identified several key pathways whose activity cycles over the circadian cycle in each cell type. As has been previously described~\cite{ref41}, we found many features of aortic metabolism undergo circadian rhythms. SMCs display rhythms in steroid biosynthesis, taurine metabolism, cholesterol metabolism, arachidonic acid metabolism, and glycolysis and gluconeogenesis. Fibroblasts display rhythms in steroid biosynthesis and fatty acid degradation. PPAR signaling peaks around ZT0 in fibroblasts and ECs. Moreover, gene expression related to cholesterol metabolism peaks around ZT0 in macrophages. As has been previously described~\cite{ref41,ref42,ref43}, SMCs demonstrate rhythmic gene expression of genes involved in contraction. Gap junction genes in SMCs peak during the dark phase, suggesting the possibility of rhythmic communication between SMCs and ECs. As has been previously described~\cite{ref44}, we found platelet activation oscillates throughout the day and peaks at ZT0 in ECs. Pathways previously described to be involved in conveying circadian time were also identified in ECs. Gene expression related to Akt signaling, which has previously been shown to signal circadian time to the vasculature~\cite{ref45}, peaked at ZT0 in ECs. Moreover, expression related to the HIF-1 pathway peaked at ZT0 in ECs, which coincides with the nadir in plasma oxygen levels~\cite{ref46}. While hypoxia has been shown to signal clock timing in mammalian contexts~\cite{ref46,ref47}, to our knowledge its role in relaying time to the vasculature has not been previously characterized and warrants further investigation. Several aspects of immune recruitment, antigen presentation, vascular adhesion, and invasion show strong circadian rhythms. VEGF signaling peaks in fibroblasts around ZT20. Chemokine signaling peaks at ZT16 in SMCs. Complement system gene expression and antigen presentation peaks at ZT18 and ZT12, respectively, in ECs. In macrophages, cytokine-cytokine receptor interaction and IL17 signaling peak around ZT0. Among all cell types, we identified extensive rhythms in focal adhesion and cell adhesion, with most genes peaking toward the start of the light phase.

\subsection*{Sex effects on the aorta circadian transcriptome}

The number of confident circadian cycling genes detected was roughly comparable between male and female cell types (\highlight{Figure 5a}), with modestly more cyclers detected in females; the fraction of cycling genes among detectable genes showed similar trends (\highlight{Figure 5a}). Female circadian gene expression was primarily distinguished by higher amplitude rhythms than observed in males (\highlight{Figure 5b}). Moreover, among circadian cycling genes common to both males and females, cyclers tended to peak modestly earlier in females than in males, which is consistent with previous observations~\cite{ref21} (\highlight{Figure 5c}).

Sex differences in the number of detected cycling genes may be due to discrepancies in the cell counts or cell library sizes. To control for discrepancies in cell counts within cell types, we downsampled male and female cells to have the same number of cells at each time point. To control for discrepancies in cell library size within cell types, we downsampled the pseudobulk UMI counts (i.e.\ total number of UMIs detected) at each time point for each condition to have the same number of pseudobulk UMIs in males and females. Thereafter, we ran our Bayesian regression pipeline. For each major cell type (\highlight{Supplementary Figure 3a}) and SMC subtype (\highlight{Supplementary Figure 3b}), we computed the number of detected circadian cyclers at various circadian Bayes factor thresholds. After adjusting for cell count and library size discrepancies, SMCs appeared more rhythmic in females, with nearly twice as many detected cycling genes at a \highlight{log10 circadian Bayes factor of 2.} However, this observation did not hold in fibroblasts, which showed similar numbers of called circadian cyclers between the sexes at all Bayes factor thresholds. As such, sex-dependent rhythmicity may be cell-type-specific. Fibroblasts and SMCs were downsampled to the same number of cells and pseudobulk UMIs, suggesting that cell type differences in sex-dependent rhythmicity were not due to discrepancies in cell or UMI counts. We note, macrophages additionally demonstrate stronger rhythmicity in females, but caution this result is less robust, due to the low number of cells detected.

While female SMCs have many more cyclers than their male counterparts, most of these are unique cyclers. To understand their functional significance, we identified overrepresented gene sets among the unique cyclers. Unique female cyclers were enriched for genes with \textit{Sp1}, \textit{Smad4, Egr1, Tead4, Sp3, Repin1,} and E-box promoter motifs. Moreover, unique female cyclers were enriched for heat shock and unfolded protein response, regulation of cellular proliferation, protein targeting to the endoplasmic reticulum, ribosome biogenesis, and negative regulation of cell death.

One of the most striking sex differences in SMCs was in the expression of the heat shock 70 related genes \textit{Hspa1a} (a.k.a.\ \textit{Hsp72}) and \textit{Hspa1b} (a.k.a.\ \textit{Hsp70-2}), and, to a lesser extent, \textit{Hspa1l} (a.k.a.\ \textit{Hsp70-1l}). \textit{Hspa1a} and \textit{Hspa1b} demonstrate nearly flat, constitutively high expression in male SMCs (\highlight{Figure 6a}). In female SMCs, \textit{Hspa1a} and \textit{Hspa1b} are among some of the most rhythmic genes, peaking at around ZT0, with substantially lower expression at ZT12 in females relative to males. This decrease in female heat shock expression at ZT12 is more prominent in SMCs of the ascending aorta/aortic arch, which may lead to an increased susceptibility to stressors at these circadian times. Estrogen has previously been shown to decrease \textit{Hsp70} expression in skeletal muscle~\cite{ref48}. Whether estrogen is responsible for the observed loss of \textit{Hsp70} expression at ZT12 remains unknown.

Beyond the role of \textit{Hsp70}'s in stress response, \textit{Hsp70} gene and/or protein expression loss is associated with multiple features of synthetic SMC phenotypes, such as increased apoptosis~\cite{ref49}, vascular calcification~\cite{ref50} and \textit{Il6} production~\cite{ref51}. Heat shock expression has been shown to inhibit SMC proliferation and mediate endothelial nitric oxide synthase (eNOS) signaling in SMCs~\cite{ref52}, and may, perhaps, underly the previously observed sex differences in eNOS sensitivity~\cite{ref53}. In addition, heat shock expression inhibits angiotensin II and MAPK signaling~\cite{ref54}. Altogether loss of heat shock expression may predispose female SMCs to phenotypic switching and the adoption of synthetic phenotypes.

To assess if sex differences in SMC phenotypic switching predisposition exist, more broadly, we looked at gene sets and signaling pathways known to promote or inhibit phenotypic switching. Several aspects of our data suggest TGF-$\beta$ signaling, which promotes the contractile SMC phenotype~\cite{ref55}, is less active in female SMCs. The TGF-$\beta$ signaling transducer \textit{Smad1} was modestly downregulated in female SMCs, particularly at ZT6 in the ascending aorta/aortic arch. Genes with \textit{Smad3} and \textit{Smad4} promoter motifs were also downregulated in female SMCs at ZT6, while \textit{Tgfb2} expression was downregulated at ZT6 in female ECs. \textit{Atf3,} an androgen inducible CREB family transcription factor involved in modulating \textit{Ap1} signaling, glucose metabolism, and stress response~\cite{ref56,ref57}, has previously been described to promote SMC identity~\cite{ref58} and phenome-wide association studies suggest its loss promotes atherosclerotic risk~\cite{ref59}. Loss of \textit{Atf3} may drive phenotypic switching, mediated by its effects on the contractile SMC transcriptional program and its effects on glucose metabolism. While \textit{Atf3} is rhythmic in both male and female SMCs, we observe an acrophase shift and amplitude increase in females, with substantially reduced expression during the light phase in females (\highlight{Figure 6a}). This decrease in female SMC \textit{Atf3} expression was particularly pronounced at ZT6 in SMCs of the ascending aorta/aortic arch. \textit{Nr4a1}, which has previously been shown to inhibit SMC proliferation and neointimal formation~\cite{ref60}, cycled in both males and females, and demonstrated decreased expression in females at ZT6 in the ascending aorta/aortic arch (\highlight{Figure 6a}). \textit{Npr3}, a receptor mediating endothelial-derived vasodilation via C-type natriuretic peptide~\cite{ref61}, cycled strongly in females and was notably downregulated at ZT6 in the ascending aorta/aortic arch (\highlight{Figure 6a}). While female SMCs may be more prone to phenotypic switching in the ascending aorta, several factors suggest they may be protected against atherosclerosis in the descending aorta. \textit{Flbn5,} which has been shown to prevent aneurysm formation and SMC proliferation and migration~\cite{ref62,ref63}, was considerably higher in female SMCs during the light phase. Moreover, our data suggested sex-specific activation of the insulin growth factor signaling (IGF) pathway~\cite{ref64} in the descending aorta, which is believed to be cardioprotective. \textit{Ghr} was flat and lowly expressed at similar levels in the ascending aorta/aortic arch of both sexes. However, \textit{Ghr} cycles strongly and is more highly expressed at ZT0/6 only in the descending aorta of females (\highlight{Figure 6a}). \textit{Igf1} was also more highly expressed at ZT0/6 of the descending aorta in females, suggesting female-specific atheroprotection at these times. Canonical contractile markers were also downregulated in females. The SMC master regulator \textit{Myocd} and genes with \textit{Myod1} promoter motifs were downregulated at ZT6 in female SMCs, a finding that was more pronounced in the ascending aorta/aortic arch. Previous work~\cite{ref65} assessing sex-dependent expression of vascular regulatory networks suggests female SMCs may have a higher predisposition for SMC phenotypic switching. While many expression changes in our data support this, a few aspects of our data conflict. \textit{Smad6} and \textit{Smad7}, which have been shown to inhibit TGF-$\beta$ signaling~\cite{ref66,ref67}, are downregulated at ZT6 in female SMCs. \textit{Nr4a3,} which is induced by PDGF~\cite{ref68}, a proatherogenic stimulus, was downregulated at ZT6. In addition, \textit{Hk2} expression was lower at ZT0/6 in female SMCs, particular in the ascending aorta, suggesting less glycolytic metabolism, which has been shown to promote phenotypic switching~\cite{ref53}. Though circadian gating of cell fate decisions has been observed in numerous contexts~\cite{ref69,ref70,ref71}, it remains unknown if and/or when the clock gates SMC phenotypic switching, and how this depends on sex. Further investigation of this matter and whether sex differences in phenotypic switching disposition are mediated by sex differences in circadian expressed should be pursued.

Beyond sex differences in the predisposition to SMC phenotypic switching, the expression profiles of female SMCs suggest they are more robust to vascular stiffness and calcification. \textit{Grem2}, a bone morphogenetic protein (BMP) antagonist~\cite{ref72}, is more highly expressed in female SMCs at ZT0/6, while genes with \textit{Sp1}~\cite{ref73} promoter motifs were enriched for downregulated female genes at ZT6.

Similar to males, females show considerable rhythms in focal adhesion related gene sets expressed in ECs, with most genes peaking around ZT0. Nonetheless, focal adhesion genes are among those enriched for downregulated mesors in females.

The data suggest various aspects of macrophage function differ between male and female mice. Male and female macrophages exhibit striking differences in \textit{Lpl}. \textit{Lpl} is constitutively highly expressed in females, with flat expression. In males, \textit{Lpl} shows strong circadian rhythms, peaking at ZT0, with substantially lower expression at ZT12 (\highlight{Figure 6b}). This may lead to sex and time dependent differences in macrophage foam cell formation and energetic metabolism. Males demonstrate stronger cycling of \textit{Ccr2,} and a modest phase delay, suggesting sex differences in the timing and amplitude of macrophage recruitment to the vasculature (\highlight{Figure 6b}). Similar to SMCs, macrophages show considerable sex differences in circadian expression of \textit{Hspa1a} and \textit{Hspa1b}, with constitutively high, flat expression males, and strong cycling with lower expression in females at ZT12 (\highlight{Figure 6b}). As \textit{Hsp70} proteins have been implicated in repressing NF-$\kappa$B activity~\cite{ref74}, we next assessed if there existed sex differences in circadian expression of classically activated macrophage markers. This analysis did not suggest clear patterns in sex-dependent expression of classical M1 vs.\ M2 macrophage gene sets, but instead suggested a more nuanced pattern of sex-dependent expression. \textit{Ccl12}, a pro-inflammatory chemokine cycled strongly in female macrophages and was upregulated at ZT0 (\highlight{Figure 6b}). \textit{Ccn2} (a.k.a.\ CTGF), which has been implicated in macrophage tissue invasion~\cite{ref75}, is decreased in expression at ZT0 in females. Genes involved in endocytosis and Hif1 signaling were enriched for downregulated genes in females at ZT12 and ZT0, respectively. Lastly, genes with \textit{Egr1} motifs were downregulated in female macrophages at ZT12, many of these genes cycling with peak times near ZT0. \textit{Egr1} has previously been shown to blunt the inflammatory response in macrophages~\cite{ref76}. While \textit{Egr1} expression was well detected in our dataset among both male and female macrophages, we did not observe sex-dependent expression, suggesting downstream mechanisms may be responsible for this observation.

\subsection*{\textit{Bmal1} knockout effects on the male aorta circadian transcriptome}

As anticipated, male \textit{Bmal1}$^{\text{fl/fl}}$:CAGG$^{\text{creERt2/+}}$ mice showed decreased rhythmic circadian transcription compared to C57BL/6J mice. Whereas hundreds of genes show strong circadian cycling in male C57BL/6J aortic cell types, only tens of such cyclers were detected in \textit{Bmal1}$^{\text{fl/fl}}$:CAGG$^{\text{creERt2/+}}$ mice (\highlight{Figure 7a}). Genes cycling in either condition show substantially reduced amplitudes in \textit{Bmal1}$^{\text{fl/fl}}$:CAGG$^{\text{creERt2/+}}$ mice (\highlight{Figure 7b}). Moreover, as expected, core clock gene expression waveforms are nearly flat (\highlight{Figure 7c}).

Several aspects suggested inducible \textit{Bmal1} KO affected genes implicated in synthetic phenotypes in SMCs. The master contractile regulator \textit{Myocd} was downregulated at ZT6, while contractile genes \textit{Myh11, Tpm1}, and \textit{Smtn} were downregulated at all time points (\highlight{Figure 8a}). This was echoed by an enrichment for genes involved in smooth muscle cell contraction among genes with decreased circadian mesors. \textit{Klf4}, whose activity has been documented to promote SMC phenotypic switching and repress myocardin-dependent contractile expression~\cite{ref39}, was upregulated at all time points but ZT6, the effects of which were more modest in the descending thoracic aorta. Genes involved in apoptosis, p53 signaling, and proteoglycans in cancer were enriched for increased mesor. Several factors in the TGF-$\beta$ pathway activation were perturbed, with \textit{Tgfbr1} downregulated at all time points, \textit{Tgfbr2} downregulated at ZT12 (\highlight{Figure 8b}), and genes with Smad4 promoter motifs enriched for downregulated genes, particularly at ZT12. However, effects on TGF-$\beta$ activation were nuanced, as inhibitors of TGF-$\beta$ signaling \textit{Smad6, Smad7, Smad9,} and \textit{Lefty1} were all downregulated at ZT12. Moreover, TGF-$\beta$ regulators \textit{Bmp4} and \textit{Bmp6} were downregulated at ZT0, while \textit{Bmp3} was down at all time points and lost cycling. Loss of \textit{Bmp} has previously been observed in phenotypic switching~\cite{ref77}. Functional assessments of alterations in TGF-$\beta$ signaling activation are warranted. Genes involved in focal adhesion, adherens junctions, and tight junctions were notably enriched among genes with downregulated mesors, suggesting SMCs may lose sources of mechanical stimuli (from both extracellular matrix and adjacent cells) essential for maintaining the contractile state~\cite{ref78,ref79}. \textit{Hk2} lost cycling and its mesor was decreased, indicating decreased glycolytic metabolism. Moreover, genes involved in choline metabolism were enriched for downregulation at ZT12, which has previously been implicated in inhibiting SMC phenotypic switching~\cite{ref80}.

In addition to alterations in contractile, TGF-$\beta$, mechanical sensing, and metabolic expression, proliferation increases are common to synthetic SMC phenotypes. However, our data suggest inducible global \textit{Bmal1} KO may repress SMC proliferation. \textit{Wee1}, a clock-controlled inhibitor of the cell cycle G2 to M transition, oscillates in all cell types and peaks around ZT12 in Bmal1$^{+/+}$ mice. In \textit{Bmal1}$^{\text{fl/fl}}$:CAGG$^{\text{creERt2/+}}$ male mice, \textit{Wee1} expression losses oscillation and is constitutively lowly expressed (\highlight{Figure 8c}). In contrast, \textit{Cdkn1a}, a clock-controlled regulator of the cell cycle G1 arrest in response to DNA damage~\cite{ref81}, oscillates in all cell types and peaks around ZT0 in Bmal1$^{+/+}$ mice. In \textit{Bmal1}$^{\text{fl/fl}}$:CAGG$^{\text{creERt2/+}}$ mice, \textit{Cdkn1a} expression loses oscillation and is constitutively highly expressed (\highlight{Figure 8c}). Existing literature~\cite{ref82} suggests \textit{Cdkn1a} activation can inhibit S-phase DNA damage induced by \textit{Wee1} inhibition. As such, \textit{Cdkn1a} expression may increase in response to \textit{Wee1} expression loss in \textit{Bmal1}$^{\text{fl/fl}}$:CAGG$^{\text{creERt2/+}}$ mice to preserve DNA integrity. However, while \textit{Cdkn1a} activity is presumed to inhibit cell cycle progression, its effects may be context dependent. Previous groups demonstrate \textit{Cdkn1a} overexpression stimulates increased proliferation in SMCs and apoptosis inhibition~\cite{ref83}. Altogether, while \textit{Wee1} and \textit{Cdkn1a} effects suggest altered SMC proliferation after \textit{Bmal1} KO, the directionality is unclear and should be the subject of future investigation.

Beyond effects on genes relevant to SMC phenotypic switching, inducible \textit{Bmal1} KO has a profound effect on sex hormone dependent gene expression. In particular, we observed several similarities between male \textit{Bmal1} KO SMCs and SMCs of female WT mice. Similar to female \textit{Bmal1} WT SMCs, \textit{Bmal1} KO SMCs downregulate \textit{Hspa1a} and \textit{Hspa1b}, particularly at ZT0 and ZT6 (\highlight{Figure 8d}). \textit{Atf3, Nr4a1,} and \textit{Nr4a3} were all downregulated at ZT6 and cycled strongly in male KO SMCs relative to WT. Previous observations made in mouse liver suggest conventional global \textit{Bmal1} KO leads to ``feminized'' expression and increased estrogen receptor activity~\cite{ref84}. Our data also reflect this finding. While moderately expressed in the descending aorta and lowly expressed in the ascending aorta with flat expression in WT mice, \textit{Esr1} is highly upregulated at all time points after inducible \textit{Bmal1} KO (though notably more highly expressed in the descending aorta). This suggests clock mechanisms may repress its expression and may mediate the cardioprotective effects observed by \textit{Yang et al.}~\cite{ref85}.

We observed large alterations in the expression of matrix metalloproteases and metalloprotease inhibitors in SMCs. \textit{Mmp14} was increased at ZT12/18 in the ascending aorta/arch, while \textit{Mmp2} was increased at all time points in all regions of the aorta. While this may suggest increased rates of extracellular matrix degradation, \textit{Timp4}, an inhibitor of metalloprotease activity, was strongly upregulated at all time points, particularly in the ascending aorta/arch.

Several effects were common to multiple cell types. In particular, decreases in focal adhesion were common to SMCs, fibroblasts, and ECs, suggesting clock function in the endogenous vascular cell types may be essential for proper mechanical sensing of extracellular matrix. Aldosterone synthesis was downregulated in SMCs and fibroblasts. Complement activation was upregulated in SMCs and fibroblasts, which has been suggested to be atherogenic~\cite{ref87,ref88}. \textit{C3, C1s1,} and \textit{C4b} were upregulated at all time points in fibroblasts (\highlight{Figure 8e}), while \textit{C1s1} and \textit{C4b} were upregulated at all time points in SMCs. \textit{C3} notably gained cycling, with increased expression at ZT6 and decreased expression at ZT18, relative to WT SMCs. This gain of SMC \textit{C3} cycling may be due to gain of cycling in \textit{Atf3}, which has been previously described to inhibit \textit{C3} transcription and cycles in antiphase in KO SMCs~\cite{ref59}.

Our results suggest \textit{Bmal1} KO may affect macrophage recruitment. While myeloid cells are rhythmically recruited to the vasculature via \textit{Ccl2-Ccr2} signaling~\cite{ref4}, with \textit{Ccr2} macrophage expression peaking at ZT0 in C57BL/6J mice, rhythmic expression of \textit{Ccr2} was delayed nearly 12 hours and its mesor was modestly increased in \textit{Bmal1}$^{\text{fl/fl}}$:CAGG$^{\text{creERt2/+}}$ macrophages (\highlight{Figure 8f}). Moreover, rhythmic signaling via the \textit{Cx3cl1-Cx3cr1} axis~\cite{ref89,ref90,ref91} appeared disrupted. \textit{Cx3cl1} and \textit{Cx3cr1} expression cycled and peaked around ZT0 in WT ECs and macrophages, respectively (\highlight{Figures 8f and 8g}). After \textit{Bmal1} KO, both genes were flat and constitutively highly expressed. Altogether, this suggests loss of \textit{Bmal1} expression may drive chronic macrophage recruitment to the aorta.

\textit{Bmal1} KO may promote a more proinflammatory environment. In ECs, expression of the anti-inflammatory signaling gene \textit{Cxcl12} was greatly decreased at ZT0, and expression of pro-inflammatory signaling gene \textit{Cd74} was increased at ZT12. In macrophages, expression of pro-inflammatory genes \textit{Cd72} and \textit{Il1b} were greatly increased at ZT18, while expression of anti-inflammatory genes \textit{Cd200} and \textit{Ccl24} was greatly decreased. While \textit{Mmp9} cycles strongly in WT male macrophages, peaking at ZT6, it is constitutively highly expressed after \textit{Bmal1} KO, suggesting KO macrophages degrade extracellular matrix and promote angiogenesis. This is likely due to the loss of \textit{Nr1d1} expression~\cite{ref92}. Genes increasing in mesor in macrophages were further enriched for positive regulation of NF$\kappa$b activity. Altogether, these results suggest that the \textit{Bmal1} KO in male mice may bias macrophages to express genes associated with a proinflammatory state, as has been previously described~\cite{ref92}.

\subsection*{Effect of sex on \textit{Bmal1} knockout effect}

KO effects in cell types are strongly concordant between male and female cell types (\highlight{Supplementary Figure 4}). In both sexes, the number of cell-type-specific confident circadian cyclers was decreased. However, females retained higher numbers of rhythmic gene sets (\highlight{Figure 9a}).

Among unique confident cyclers in female SMCs, dusk cyclers were enriched for TNF-$\alpha$ signaling and inhibition of MAPK signaling, while dawn cyclers were enriched for ECM organization and Pi3k-Akt signaling. Moreover, estrogen signaling genes were enriched at dusk, particularly in SMCs of the ascending aorta/aortic arch. Similar to its expression in male SMCs, \textit{Esr1}'s mesor greatly increases in female SMCs after \textit{Bmal1} KO. Unlike in males, \textit{Esr1} cycles strongly in female SMCs, peaking at ZT15 (\highlight{Figure 10a}). Several other components of the estrogen signaling pathway uniquely cycled in female KOs, such as \textit{Hsp90aa1, Fkbp5, Hbegf, Atf4}, peaking at dusk. Among estrogen signaling cycling genes induced by KO were \textit{Jun} and \textit{Fos}, which gained expression at ZT12, specifically in the ascending aorta/aortic arch of female SMCs. This suggests increased \textit{Ap1} transcription factor activity at ZT12, which has been implicated in SMC proliferation and neointimal formation~\cite{ref93,ref94,ref95}. Interestingly, \textit{Bmal1} KO in female SMCs induces an acrophase shift of \textit{Atf3} to ZT12. As \textit{Atf3} may antagonize \textit{Ap1} activity~\cite{ref96}, this may reflect a compensatory effect to buffer against changes in \textit{Ap1} activity after KO.

One of the most striking sex differences in \textit{Bmal1} WT SMCs was in the expression of \textit{Hspa1a} and \textit{Hspa1b} (\highlight{Figure 10a}). Both were among the strongest cyclers in female SMCs, peaking around ZT21, with considerably lower expression in females at ZT12. However, after \textit{Bmal1} KO, \textit{Hspa1a} and \textit{Hspa1b} appeared similar between the sexes, cycling in both sexes with an acrophase around ZT15 and modestly lower expression in females during the light phase. Previous reports suggest \textit{Bmal1} KO leads to ``feminized'' circadian expression liver, and that these effects are likely due to loss of central oscillators~\cite{ref84}. Our data suggest expression of WT confident cyclers in SMCs becomes more similar between sexes after KO (\highlight{Figures 10b and 10c}). This suggests that sex-specific expression may be mediated, in large part, through expression of circadian cycling genes.

Several aspects suggest female SMCs may proliferate less than male SMCs after \textit{Bmal1} KO. \textit{Mcm6}, a mini-chromosome maintenance protein required for replication, cycles in female WT SMCs, with an acrophase around ZT15. In female SMCs, \textit{Mcm6} expression is downregulated at all time points, an effect not observed in males (\highlight{Figure 10a}). \textit{Plat}, a plasminogen activator that has previously been shown to be a mitogen for SMCs~\cite{ref97}, loses expression uniquely in females at ZT12 and ZT18 (\highlight{Figure 10a}). This effect is more pronounced in SMCs of the ascending aorta/aortic arch. Similarly, \textit{Apoe,} which can inhibit SMC proliferation by activating inducible nitric oxide synthase~\cite{ref98}, gains expression more strongly in females at ZT12 and ZT18 than in males (\highlight{Figure 10a}), which is particularly pronounced in the ascending aorta/aortic arch. This same trend was observed for \textit{Errfi1}, which inhibits neointimal formation~\cite{ref99}. Nonetheless, \textit{Rgs5}, which inhibits SMC proliferation~\cite{ref31,ref100,ref101} by antagonizing G-coupled protein receptor signaling, demonstrates decreased expression at ZT6 in the ascending aorta/aortic arch and at all time points in the descending thoracic aorta. In the absence of G-coupled protein receptor signaling, the net effect may still be an inhibition of SMC proliferation in female KO SMCs. Nonetheless, understanding how these effects on key proliferation genes ultimately affects proliferation will require experimental follow up.

Female SMCs may also be less likely to migrate after \textit{Bmal1} KO. \textit{Serpine1} (a.k.a.\ \textit{Pai1}), which can stimulate SMC migration and neointimal formation~\cite{ref102}, is downregulated at ZT0 in the ascending aorta/aortic arch of females relative to their male KO counterparts (\highlight{Figure 11a}); \textit{Ace}, an angiotensin converting enzyme promoting SMC migration~\cite{ref103}, is only downregulated in female SMCs of the aortic arch/ascending aorta at ZT12/18 (\highlight{Figure 11a}). Female SMCs may also be more predisposed to oxidative stress after KO. \textit{Nupr1}, a key gene in antioxidant activity~\cite{ref104}, is specifically downregulated in female SMCs of the aortic arch/ascending aorta at ZT12/18.

While globally upregulated in male SMCs after \textit{Bmal1} KO, \textit{Cdkn1a} expression was not altered by \textit{Bmal1} KO in male macrophages. However, macrophage \textit{Cdkn1a} expression was strongly downregulated at ZT6 in female macrophages, suggesting an increased likelihood to progress through the cell cycle at these time points (\highlight{Figure 11b}). Several KO effects observed in male macrophages, were shared by females, such as an increased mesor for \textit{Cd72} and \textit{Il1b}. However, both genes displayed acrophase shifts. Similarly, \textit{Ccr2} peaked around ZT15 in male macrophages after \textit{Bmal1} KO, while \textit{Ccr2} peaks around ZT20 in female macrophages, suggesting sex dependent differences in myeloid recruitment after clock KO (\highlight{Figure 11b}). Genes with mesor decreases in female KO relative male KO macrophages were enriched for lipopolysaccharide response genes and cellular response to cytokine stimulus, such as \textit{Zfp36, Ccl3, Tnfaip3, Ccl2, Cd14, Tnf}, and \textit{Ccl4}, with lower expression in females KO macrophages around ZT6--ZT12. Moreover, genes downregulated at ZT6 in female KO macrophages relative to male KO macrophages were enriched for NF-$\kappa$B signaling. Altogether, these results support that \textit{Bmal1} KO biases female macrophages toward an anti-inflammatory phenotype.

\subsection*{Misalignment effects on the male aorta circadian transcriptome}

To profile the circadian transcriptomes of mice acutely misaligned with the LD cycle, male C57BL/6J mice were subjected to a 6-hour phase advance and aorta transcriptomes were profiled 5 days later at ZT0, ZT6, ZT12, and ZT18 in the new LD cycle. The number of confident circadian cycling genes was decreased after misalignment (\highlight{Figure 12a}). While this may be due to misalignment, this may also have been explained by fewer cells captured in the misaligned condition, and therefore less power for cycling detection. Nonetheless, after adjusting for discrepancies for cell counts and library sizes, fewer cycling genes were still detected within each cell type after misalignment, and this was robust to the Bayes factor threshold used to called cyclers (\highlight{Supplementary Figure 5}). For confident cycling genes common to the aligned and misaligned condition within each cell type, cycling gene acrophases were generally delayed in the misaligned mice (\highlight{Figure 12b}), indicating the mice had not completely entrained to their new schedule. In addition, amplitudes of these confident cycling genes were generally decreased in the misaligned condition (\highlight{Figure 12c}).

In male misaligned SMCs, heat-shock protein expression is strikingly altered. While in the aligned male SMCs, \textit{Hspa1a} and \textit{Hspa1b} are highly expressed and with modest circadian cycling, \textit{Hspa1a} and \textit{Hspa1b} is greatly decreased at ZT0 in male misaligned mice (\highlight{Figure 13}). Other heat shocked related genes show similar patterns, such as \textit{Dnajb1, Dnajb4, Dnaja1,} and \textit{Hsph1}, suggesting male SMCs may be more susceptible to stressors around ZT0 after acute misalignment. Beyond their roles in heat shock response, heat shock proteins are also implicated in endoplasmic reticulum (ER) stress and the unfolded protein response. Additional aspects of our data suggest SMCs may be susceptible to these stressors. \textit{Xbp1,} a key regulator of the unfolded protein response and which has been implicated in preventing neointimal hyperplasia~\cite{ref105}, is downregulated at ZT0 (\highlight{Figure 13}). \textit{Ankrd1}, which protects cells from ER stress~\cite{ref106}, is downregulated during the light phase in the descending aorta (\highlight{Figure 13}). \textit{Paip2b}, which inhibits mRNA translation~\cite{ref107}, is downregulated at all time points, which may lead to the accumulation of unfolded proteins (\highlight{Figure 13}). We additionally observe other alterations in gene expression that may aid in reducing the number of unfolded proteins. \textit{Uba52,} which targets proteins for degradation via its ubiquitin function and has effects on cell cycle and ribosomal regulation~\cite{ref108,ref109}, is strongly upregulated at all time points (\highlight{Figure 13}). Genes with decreased mesors are enriched for transcription and ribosome biogenesis, which may reflect an effort to decrease transcription. Beyond ER stressors, SMCs may also be more susceptible to DNA mutagens, as \textit{Wdr70,} a key factor in processing double strand DNA breaks~\cite{ref110,ref111}, is downregulated at all time points after misalignment (\highlight{Figure 13}).

We additionally observe broad alterations in gap junction expression, the effects of which appear more pronounced in SMCs of the ascending aorta/aortic arch relative to the descending aorta. Several tubulin genes involved in \textit{Cx43} trafficking are downregulated over the circadian cycle, and \textit{Cx43} expression is notably decreased at ZT0 in SMCs, which is mirrored in ECs. These results point to dysregulated SMC-endothelial signaling upon misalignment with the LD cycle.

Male misaligned SMCs exhibit expression changes associated with multiple features of synthetic phenotypes. \textit{Atf3}, a transcription factor involved in cAMP signaling and stress response, has previously been described to inhibit SMC phenotypic switching. The \textit{Atf3} acrophase is delayed by around 7 hours and expression is substantially decreased in male misaligned mice at ZT0 (\highlight{Figure 13}), suggesting acute misalignment may predispose SMCs to phenotypic switching at the start of the circadian cycle. \textit{Notch2}, which has been previously demonstrated to inhibit SMC phenotypic switching and Pdgf-B-dependent proliferation~\cite{ref112,ref113,ref114} was downregulated. Meanwhile, \textit{Rock1} and \textit{Egln1}, which both promote Pdgf-B-dependent proliferation~\cite{ref115,ref116}, had increased mesors. As the descending thoracic aorta is generally more atheresistant, it may also be more resistant to perturbations that induce phenotype switching. We noted \textit{Ccn1}, the loss of which has been shown to prevent neointimal hyperplasia~\cite{ref117}, is downregulated during the dark phase only in SMCs of the descending aorta.

Our results suggest SMC energetic metabolism may be altered after misalignment. In particular, male misaligned SMCs downregulate \textit{Hk2} at ZT0, suggesting decreased glycolytic metabolism (\highlight{Figure 13}). This contrasts with the general finding that synthetic SMCs increasingly rely on glycolytic metabolism and fatty acid oxidation instead of the citric acid cycle~\cite{ref53}. This may reflect an effort by SMCs to maintain a canonical contractile state after being perturbed via misalignment.

In ECs, acute misalignment leads to reduced expression of inhibitors of angiogenesis and VEGF signaling, such as \textit{Rgs5} and \textit{Ogn} (\highlight{Figure 14a}), and reduced expression of \textit{Igfbp5}, which inhibits endothelial proliferation (\highlight{Figure 14a}). This suggests acute misalignment may alter endothelial barrier function, which could make the vasculature more susceptible to immune infiltration. Endothelial TGF-$\beta$ signaling expression is decreased at ZT0, particularly for \textit{Tgfb1} (\highlight{Figure 14a}) and \textit{Tgfb3}, which may additionally predispose SMCs to proliferation, migration, and phenotypic switching.

Acute misalignment in male mice may alter macrophage recruitment and polarization. While most common cycling genes between aligned and misaligned conditions had expected acrophase differences of less than 6 hours, suggesting active entrainment to the new LD cycle, both \textit{Ccr2} and \textit{Cx3cr1} had expected acrophase differences of around 6 hours (\highlight{Figure 14b}). This suggests that \textit{Ccr2} and \textit{Cx3cr1} have changed their acrophase peak times very little relative to the old LD cycle, and may lead to aberrantly timed recruitment of macrophages after acute misalignment.

To assess if there exist cell-type-specific differences in entrainment speed, we assessed acrophase entrainment of the core clock genes in SMCs, fibroblasts, ECs, and macrophages. We limited this analysis to clock genes that confidently cycled in all major cell types of both male aligned and misaligned mice. For each gene in each cell type, we then estimated the posterior distribution of the acrophase delay in the misaligned condition relative to the aligned condition. If genes are perfectly entrained to the new LD cycle after misalignment, the acrophase delay should be 0 hours. Likewise, if genes did not adapt their timing to the new LD cycle, the acrophase delay should be 6 hours (given that the misaligned mice were phase advanced 6 hours). For each gene in each cell type, we visualized the acrophase delay posterior distributions (\highlight{Figure 15}) and computed the fraction of samples between 4.5 and 7.5 hours (i.e.\ a 3 hour interval centered at 6 hours) as a measure of posterior likelihood that the gene did not shift. In general, for clock genes, the point estimate of the posterior acrophase delay was larger in macrophages than in other cell types. Moreover, for all clock genes evaluated, the posterior acrophase delay distributions of macrophages demonstrated the most posterior density near 6 hours, suggesting macrophage clocks may shift less quickly than clocks of the endogenous vascular cell types. A potential complication in the interpretation of these results is the large difference in cell counts between the cell types. Nonetheless, ECs and macrophages have comparable cell counts (466 and 423 cells in aligned male ECs and macrophages; 265 and 217 cells in misaligned ECs and macrophages). Across all clock genes evaluated, the posterior acrophase delay demonstrated more posterior density near 6 hours in macrophages than ECs. Altogether, this suggests misalignment may lead to a disruption of the relative circadian timing between cell types after acute misalignment.

\subsection*{Sex effects on misalignment}

Similar to the observations made in the aligned mice, after acute misalignment, females had more confident cycling genes (\highlight{Figure 16a}). This finding held after adjusting for discrepancies in cell counts and library sizes (\highlight{Supplementary Figures 5 and 6}). Confident cycling genes common to both males and females demonstrate larger amplitudes (\highlight{Figure 16b}). We next evaluated whether female C57BL/6J cell types more quickly entrain to their new LD cycle. To inspect this, we restricted our analysis to confident cycling genes common to the male aligned, male misaligned, female aligned, and female misaligned conditions within SMC subtypes and fibroblasts. For each gene in each sex and cell type, we computed the posterior distribution of the absolute value of the difference between the misaligned and aligned peak times (i.e.\ acrophase difference distributions). Smaller values indicate better entrainment (i.e.\ after misalignment, gene exhibits more similar peak time to the aligned peak time). We then asked if we sampled from both the male and female acrophase difference distributions, what fraction of the time would the samples be smaller in females. Females demonstrated a consistent, albeit modest, improvement in cell type entrainment (\highlight{Supplementary Figure 7}) for the SMC subtypes and fibroblasts.

After acute misalignment with the LD cycle, female SMCs demonstrate transcriptomic changes similar to male SMCs, regardless of expression differences in the aligned condition. For example, \textit{Uba52} and \textit{Mcm6} cycle in aligned female SMCs and weak cycling in male SMCs, with downregulated expression in aligned female SMCs at ZT6 relative to male SMCs. However, after misalignment, \textit{Uba52} and \textit{Mcm6} lose cycling and gain mesor, with expression at similar levels in male and female SMCs.

Several factors implicated in SMC phenotypic switching are altered specifically in female SMCs after acute misalignment. \textit{Id1} and \textit{Id3} expression is lost at ZT12 in all aortic regions in female SMCs, with minimal expression changes in male SMCs (\highlight{Figure 17a}). In SMCs, \textit{Id1} and \textit{Id3} are induced by mitogens, regulate SMC proliferation, and may inhibit differentiation~\cite{ref118,ref119,ref120}, suggesting their female-specific perturbation after misalignment may affect synthetic aspects of SMC function. In alignment \textit{Ccn2} -- which upregulates \textit{Tgfbr2}~\cite{ref121}, regulates adhesion, and mediates heparin signaling~\cite{ref122}, which inhibits SMC proliferation~\cite{ref53} -- cycles strongly, peaking around ZT21 in male and female SMCs. After misalignment, \textit{Ccn2} loses cycling more strongly in female SMCs, with considerably less expression in misaligned female SMCs relative to aligned female SMCs at ZT18, particularly in the ascending aorta/aortic arch (\highlight{Figure 17b}).

Several aspects of our data suggest endothelial barrier function is decreased in males, but not females, after acute misalignment with the LD cycle. \textit{Egr3}, which is induced after VEGF signaling~\cite{ref123,ref124}, is downregulated at ZT0/18 in female misaligned ECs relative to male misaligned ECs (\highlight{Figure 18a}). \textit{Nrarp}, which is implicated in angiogenesis~\cite{ref125}, was down at ZT6 in misaligned females relative to misaligned males (\highlight{Figure 18a}). \textit{Fabp4}, inhibition of which decreases \textit{VEGF} signaling and angiogenesis~\cite{ref126,ref127}, was downregulated at ZT0 in misaligned females relative to misaligned males (\highlight{Figure 18a}). Altogether, this suggests males may have decreased endothelial barrier after acute misalignment.

After misalignment, female macrophages lose \textit{Dusp2} expression at ZT6/12, a downstream target of NF-$\kappa$b~\cite{ref128} and the loss of which has been shown to reduce macrophage secretion of proinflammatory mediators~\cite{ref129}. After misalignment, \textit{Hbegf} expression is lost only in female macrophages. \textit{Hbegf} activity has previously been implicated in promoting macrophage-driven angiogenesis~\cite{ref130,ref131} and macrophage-driven SMC proliferation~\cite{ref132}. Moreover, after misalignment, \textit{S100a6} expression is only gained in female macrophages at ZT0. Previous observations suggest \textit{S100a6} activates RAGE signaling~\cite{ref133}, which may drive an increase in female macrophage apoptosis at ZT0.

%-----------------------------------------------------------------------
% DISCUSSION
%-----------------------------------------------------------------------
\section*{Discussion}

Circadian clock disruption has a notable influence of cardiovascular health. Defining circadian clock outputs in health are a key first step to understanding how clock perturbations mediate downstream cardiovascular effects. Bulk RNA-sequencing circadian time courses of the aorta have been a boon to defining key vascular functions cycling with circadian rhythms. Nonetheless, the cell-type-specific contribution to circadian clock outputs, and how they differ by sex, remains unexplained. In this study, we generate an atlas of cell-type-specific aortic circadian expression in both male and female mice using scRNA-seq. These data suggest $\sim$5--10\% of the cell-type-specific transcription confidently cycles with a circadian rhythm, in both males and females.

We identify several novel sex-dependent aspects of circadian expression within cell types. Within SMCs, we identify strong sex-dependent patterns of heat shock gene expression. More specifically, we identify substantially decreased expression of \textit{Hspa1a} and \textit{Hspa1b} at ZT12. This may mediate sex-dependent and time-dependent differences in the response to cellular stressors in SMCs. Beyond their roles in heat shock response, \textit{Hspa1a} and \textit{Hspa1b} expression loss may mediate the promotion of SMC synthetic phenotypes. Assessing sex-dependent expression effects on gene sets mediating synthetic phenotypes suggested (1) females may be susceptible to phenotypic switching, more broadly (2) that sensitivity to atherogenic stimuli may depend on time of day and (3) that these effects are more pronounced in the atheroprone ascending aorta/aortic arch region. Future functional studies are needed to detail whether SMC phenotypic switching is gated by the circadian cycle, and whether the sensitivity to stimuli promoting phenotypic switching differs between the sexes.

We next probed how cell-type-specific transcriptomes are perturbed by inducible deletion of \textit{Bmal1}. In male SMCs, loss of \textit{Bmal1} increases \textit{Esr1} expression and promotes more ``feminized'' expression patterns. \textit{Bmal1} loss may promote synthetic phenotypes, as contractile gene expression is lost and \textit{Klf4} is globally upregulated. Across all cell types, \textit{Bmal1} loss affected gene sets involved in cellular proliferation and focal adhesion, suggesting these functions are under circadian clock control. Previous observations suggest loss of endogenous vascular clock function promotes macrophage and T cell infiltration~\cite{ref13}. Our data suggests that this may be mediated by altered \textit{Cx3cl1 -- Cx3cr1} and \textit{Ccl2 -- Ccr2} signaling between ECs and macrophages after \textit{Bmal1} deletion.

We next probed how cell-type-specific transcriptomes are perturbed by acute misalignment with the LD cycle, via a 6-hour phase advance. Our study highlights several cell-type-specific effects that may underly the deleterious cardiovascular effects of circadian misalignment with the LD cycle. SMCs may be more susceptible to phenotypic switching, due to alterations in TGF-$\beta$ gene sets both in SMCs and ECs. SMCs may undergo increased ER stress after acute misalignment by the LD cycle. Downregulation of \textit{Paip2b} suggests the number of unfolded proteins may increase in SMCs after acute misalignment. Nonetheless, SMCs may be less equipped to deal with this increased burden of unfolded proteins. Unfolded protein response genes, such as \textit{Xbp1} and many heat-shock related genes, are downregulated during the light phase, suggesting SMCs may be less able to activate the unfolded protein response after acute misalignment. We observe several expression changes that may reflect a compensatory effort by SMCs to decrease unfolded protein burden. \textit{Uba52}, which promotes protein degradation is strongly upregulated, and gene sets involved in promoting transcription and ribosome biogenesis are downregulated. Follow up experiments are required to understand whether circadian misalignment, indeed, increases the unfolded protein burden and whether the unfolded protein response is dysregulated. SMCs may also be more prone to phenotypic switching and susceptible to DNA double strand breaks after acute misalignment. After circadian misalignment, vascular barrier function may be hampered, as endothelial cell expression changes suggest increased VEGF and angiogenesis related signaling. Circadian misalignment may also perturb the relative timing between vascular cell types and macrophages. \textit{Cx3cr1} and \textit{Ccr2} rhythmic expression mediates the rhythmic recruitment of macrophages to the vasculature. While macrophage rhythms of \textit{Cx3cr1} and \textit{Ccr2} expression are preserved after acute misalignment, their timing does not adapt to the new LD cycle. As such, macrophages may be recruited at the ``wrong'' times. Assessing whether slower entrainment to the LD cycle was a general characteristic of macrophage clocks, we assessed core circadian clock gene expression entrainment to the new LD cycle in all cell types. Indeed, our analyses suggest macrophage clocks may entrain more slowly than the endogenous vascular cell types. As such, the relative timing between clocks may be altered during acute circadian misalignment, which may be deleterious. Nonetheless, our study captured relatively few macrophages. Future efforts, ideally using multi-channel \textit{Per2::luciferase} cell-type-specific reporters \textit{in vivo}, should be used to resolve whether the relative clock timing between macrophages and endogenous vascular cell types is disrupted after misalignment with LD cycle.

Lastly, we probed the sex-specific response to acute misalignment with the LD cycle. Acute misalignment exacerbated the female-specific increase in factors promoting SMC phenotypic switching. Nonetheless, these factors appear distinct from those observed after \textit{Bmal1} deletion, and are primarily defined by loss of \textit{Id1, Id3}, and \textit{Ccn2} expression. Female vasculature may be protected against the barrier function effects observed in males, as gene sets promoting VEGF and angiogenesis signaling in ECs were regulated relative to males after misalignment.

While our study highlights several novel aspects of cell-type-specific and sex-dependent circadian transcription, and how it is affected by various perturbations, it is not without pitfalls. In particular, cell-type-specific gene expression comparisons between sexes were mostly performed by directly comparing posterior expression levels of individual genes. Nonetheless, this may lead to misleading interpretations. For example, sex-dependent regulatory networks could generate the same activation of a signaling pathway given differing levels of gene expression levels within the signaling pathway. Moreover, while our study collected many SMCs, allowing us to study circadian expression with high confidence and to probe regional aortic variation in expression, we collected relatively few macrophages and ECs. As such, the number of confident circadian cyclers detected in these cell types is likely an underestimate, and peak time estimates are more uncertain. As such, future work remains to characterize the endothelial cell and macrophage circadian transcriptions with more detail.

Altogether, we generated a circadian atlas of cell-type-specific and sex-dependent expression in the aorta, and characterized how this expression is perturbed by adult-life inducible genetic deletion of \textit{Bmal1} and misalignment with the LD cycle. We hope this will provide the foundation for future studies to explore the mechanistic basis of how circadian disruption affects cardiovascular health and how these effects may differ by sex.

%-----------------------------------------------------------------------
% METHODS
%-----------------------------------------------------------------------
\section*{Methods}

\subsection*{Inducible deletion of the \textit{Bmal1} gene in \textit{Bmal1}$^{\text{fl/fl}}$:CAGG$^{\text{creERt2/+}}$ mice}

At 8 weeks of age, \textit{Bmal1} was inducibly deleted in \textit{Bmal1}$^{\text{fl/fl}}$:CAGG$^{\text{creERt2/+}}$ mice by administering tamoxifen via oral gavage (100 mg/kg of body weight) for 5 days. Confirmation of \textit{Bmal1} protein expression loss was confirmed via Western blot (\highlight{Supplementary Figure 8}).

\subsection*{Aorta collection and scRNA-seq data generation}

For the aligned conditions, 12-week-old C57BL/6J and \textit{Bmal1}$^{\text{fl/fl}}$:CAGG$^{\text{creERt2/+}}$ male and female mice from the Jackson Laboratory were entrained to a 12:12 LD cycle (12 hours of light, followed by 12 hours of darkness) for 2 weeks in circadian boxes with ventilation (20--22$^{\circ}$C, RH $\sim$50\%) and fed \textit{ad libitum}. At ZT0, ZT6, ZT12 (lights off), or ZT18 mice were sacrificed (2 per time point) via CO$_2$ exposure.

For the misaligned conditions, 12-week-old C57BL/6J male and female mice from the Jackson Laboratory were entrained to a 12:12 LD cycle (12 hours of light, followed by 12 hours of darkness) for 2 weeks in circadian boxes with ventilation (20--22$^{\circ}$C, RH $\sim$50\%) and fed \textit{ad libitum}. Mice were then phase advanced 6 hours. 5 days after the phase advance, mice were sacrificed at ZT0, ZT6, ZT12 (lights off), or ZT18 (2 per time point).

For each mouse in each condition, whole aorta was dissected and cut into small pieces ($\sim$1 mm), and cells were disassociated with 1.5 mL of enzyme cocktail (DNase--120 U/mL (Worthington, \#LS006331), Liberase TM--4 U/mL (Roche, \#05401127001), hyaluronidase--60 U/mL (Sigma-Aldrich, \#H3506)) in a petri dish at 37$^{\circ}$C for 40 mins. A pipette (1 mL) was used to dissociate aortic cells from the tissue pieces every 10 mins during the incubation. Cell supernatant was filtered through a 40 $\mu$m strainer and washed with RPMI1640 (Gibco, \#1187-085) containing 10\% fetal bovine serum (FBS, HyClone, \#SH30071.03) to inactive the enzyme cocktail. Residual red blood cells were lysed by incubating the cell suspension with Red Blood Cell Lysing Buffer Hybri-Max (Sigma-Aldrich, \#R7767) at RT for 1 min. The cells were washed two more times to remove debris with FACS buffer (FBS 2\%), EDTA (5 mM, Invitrogen, \#15575-038), HEPES (20 mM, Gibco, \#15630-080), sodium pyruvate (1 mM, Gibco, \#11360-070) in 1x PBS (Gibco, \#14190-136), and resuspended in DMEM/F12 (Gibco, \#11320-033) media containing 10\% FBS for further analysis. Cells from mice sacrificed at the same time point were pooled to form single-cell suspensions. Single cells were then captured and barcoded using the 10X Genomics Chromium platform and sequenced using an Illumina NovaSeq S2 flow cell. Cell barcode detection, read alignment, and transcript quantification were performed using the 10X Genomics Cell Ranger pipeline.

\subsection*{Cell type marker gene detection procedure}

Cell type markers used as input for low-dimensional embedding estimation were selected via a three-step procedure yielding 1969 marker genes.

First, markers whose variation defines major cell types were selected via a prior knowledge-based procedure. This was performed because preliminary clustering suggested high cell type imbalance of our dataset, with a disproportionate amount of SMCs and fibroblasts. As such, markers of rarer cell types (e.g.\ T cells) may not be identified using traditional methods based on detecting highly variable genes. Based on our prior work described in \textit{Pan et al.}~\cite{ref35}, we anticipate SMCs, fibroblasts, ECs, macrophages, and T cells. Based on this work, we selected several highly expressed and highly specific markers for each major cell type. We selected \textit{Myh11} and \textit{Acta2} for SMCs; \textit{Serpinf1} and \textit{Serping1} for fibroblasts; \textit{Vcam1} and \textit{Pecam1} for endothelial cells; \textit{Cd68, Cd14, Cd16,} and \textit{Cd64} for macrophages; and \textit{Cd4, Cd8, Cd3g,} and \textit{Cd3d} for T cells. For each major cell type, we then identified confident cells belonging to that cell type based on the transformed expression of their corresponding prior knowledge marker genes. Expression was transformed as follows:

\begin{equation}
X_{ij}^{(\text{transformed})} = \log_{10}\!\left( \min\!\left( \frac{X_{ij}}{L_{i}},\, 10^{-7} \right) \right)
\end{equation}

Where $X_{ij}$ is the UMI count of gene $j$ in cell $i$ and $L_{i}$ is the library size of cell $i$. Cells with transformed expression greater than $-3$ for any cell type marker was assigned to the corresponding cell type. This led to 119,817, 18,547, 3,632, 525, 185 cells roughly labeled as SMCs, fibroblasts, ECs, macrophages, and T cells. 5498 cells remained unlabeled.

Second, we identified markers with highly correlated expression to the prior knowledge cell type markers after the rough cell assignments. To identify these markers, the $\log_{10}$ pseudobulk relative expression of all genes was computed within each rough cell type. Cell type markers were then identified as those with $\log_{10}$ pseudobulk relative expression of at least $-4$ (i.e.\ 1 in 10,000 transcripts in pseudobulk) and with $\log_{10}$ pseudobulk relative expression larger than all other cell types by at least 0.5. This yielded 114, 113, 101, 228, 93 additional markers for SMCs, fibroblasts, ECs, macrophages, and T cells, respectively.

Third, we identified markers with high variability within cell types. To identify these marker genes, we first limited ourselves to genes that were sufficiently highly expressed. We decided cell type marker genes should be expressed at relative proportions of 1 in 10,000 transcripts in the cell types they mark. Moreover, there should be at least 30 cells observed for each cell type. Given this and given the median library size of 12,553 UMIs in our dataset, cell type marker genes should have at least 37 pseudobulk UMI for the cell type of interest (i.e.\ $10^{-4} \times 30 \times 12553$). Using this set of genes, we next restricted genes to those that had highly variable expression. Using the transformed expression counts, we computed the mean and variance of each genes' transformed expression. After fitting a nonparametric kernel (Python statsmodels implementation of Nadaraya-Watson, using bandwidth of 0.1 and default values for all other parameters) to describe the mean-variance relationship, genes with large variances were detected as those with Pearson residuals larger than 0.5.

\subsection*{Bayesian temporal regression model}

\subsubsection*{Likelihood model}

A Bayesian ANOVA regression model was fit to describe gene waveforms and detect genes with circadian rhythms. Our model was fit independently for each cluster and condition combination. As input, our model requires an $n \times p$ transcript count matrix $\mathbf{X}$, where $n$ is the number of cells and $p$ is the number of genes. For gene $j$ in cell $i$ with library size $L_{i}$ and cell phase $\Theta_{i}$ (where $\theta_{i} \in \{1,\ldots,q\}$ discrete time points), the unique molecular identifier (UMI) count $X_{ij}$ was assumed to follow a Negative Binomial (NB) distribution:

\begin{equation}
X_{ij} \mid (\Theta_{i} = k) \sim \mathrm{NB}\!\left( L_{i}\lambda_{jk},\,\delta_{ij} \right)
\end{equation}

Where $k$ is the discrete index of the time the cell was collected at,

\begin{equation}
\mathbb{E}\!\left[ X_{ij} \right] = L_{i}\lambda_{jk},
\end{equation}

\begin{equation}
\mathrm{Var}\!\left( X_{ij} \right) = \mathbb{E}\!\left[ X_{ij} \right] + \delta_{jk} \times \left( \mathbb{E}\!\left[ X_{ij} \right] \right)^{2},
\end{equation}

and

\begin{equation}
\delta_{jk} = g_{\boldsymbol{\zeta}}\!\left( \lambda_{jk} \right),
\end{equation}

and where $g_{\boldsymbol{\zeta}}\!\left( \lambda_{jk} \right)$ is a deterministic polynomial function parameterized by $\boldsymbol{\zeta}$ (shared by all cells and genes) describing the relationship between transcript proportion $\lambda_{jk}$ and the dispersion $\delta_{jk}$. Details on the estimation of $\boldsymbol{\zeta}$ can be found in \highlight{Supplementary Methods 1}.

\subsubsection*{Prior knowledge of gene parameters}

Prior knowledge about $\lambda_{jk}$ is specified as a transformed Beta distribution:

\begin{equation}
P\!\left( \lambda_{jk} \right) = \mathrm{Beta}\!\left( \lambda_{jk}^{(\alpha)},\lambda_{jk}^{(\beta)} \right) \times \left( \lambda_{j}^{(\max)} - \lambda_{j}^{(\min)} \right) + \lambda_{j}^{(\min)}
\end{equation}

where $\lambda_{j}^{(\min)}$ and $\lambda_{j}^{(\max)}$ denote the minimum and maximum possible values for $\lambda_{jk}$. $\lambda_{j}^{(\min)}$ and $\lambda_{j}^{(\max)}$ are set as follows. For each gene we compute the pseudobulk relative proportion (i.e.\ $\frac{\sum_{i=1}^{p} X_{ij}}{\sum_{i=1}^{p} L_{i}}$). Assuming relatively balanced numbers of cells from each collection time point, the pseudobulk relative proportion is a reasonable estimate of each genes' mesor. Moreover, based on analysis of past scRNA-seq circadian datasets, the largest observed fold changes of cycling genes from mesor to peak are around 10-fold (i.e.\ 1 in $\log_{10}$ space), while most have fold changes of around 2-fold. As such, we define $\lambda_{j}^{(\min)}$ and $\lambda_{j}^{(\max)}$ to account for fold changes of around 30-fold to be conservative (1.5 in $\log_{10}$ space):

\begin{equation}
\log_{10}\!\left( \lambda_{j}^{(\min)} \right) = \log_{10}\!\left( \frac{\sum_{i=1}^{p} X_{ij}}{\sum_{i=1}^{p} L_{i}} \right) - 1.5
\end{equation}

\begin{equation}
\log_{10}\!\left( \lambda_{j}^{(\max)} \right) = \log_{10}\!\left( \frac{\sum_{i=1}^{p} X_{ij}}{\sum_{i=1}^{p} L_{i}} \right) + 1.5
\end{equation}

By default, \highlight{Tempo} sets $\lambda_{jk}^{(\alpha)} = \lambda_{jk}^{(\beta)} = 1$, which assumes a non-informative prior over the domain of possible $\lambda_{jk}$ values.

\subsubsection*{Approximate posterior inference}

Using our prior knowledge of the gene parameters and the observed data, we seek the following joint posterior distribution of our gene parameters:

\begin{equation}
P\!\left( \boldsymbol{\lambda} \mid \mathbf{X}, \boldsymbol{\theta} \right) \sim P\!\left( \mathbf{X} \mid \boldsymbol{\lambda}, \boldsymbol{\theta} \right) P\!\left( \boldsymbol{\lambda} \mid \boldsymbol{\theta} \right)
\end{equation}

where $\boldsymbol{\lambda}$ is a $p \times q$ dimensional matrix containing the transcript proportion parameters for each gene at each time point, and

\begin{equation}
P\!\left( \mathbf{X} \mid \boldsymbol{\lambda}, \boldsymbol{\theta} \right) = \prod_{i=1}^{n} \prod_{j=1}^{p} \prod_{k=1}^{q} P(X_{ij} \mid (\Theta_{i} = k))
\end{equation}

and

\begin{equation}
P\!\left( \boldsymbol{\lambda} \mid \boldsymbol{\theta} \right) = \prod_{j=1}^{p} \prod_{k=1}^{q} \lambda_{jk}
\end{equation}

No known analytic solution exists to $P\!\left( \boldsymbol{\lambda} \mid \mathbf{X}, \boldsymbol{\theta} \right) \sim P\!\left( \mathbf{X} \mid \boldsymbol{\lambda}, \boldsymbol{\theta} \right) P\!\left( \boldsymbol{\lambda} \mid \boldsymbol{\theta} \right)$. Moreover, asymptotically exact estimation approaches, such as Markov-chain Monte Carlo and full grid sampling, do not scale well to droplet-based scRNA-seq datasets that can contain thousands (and sometimes tens of thousands) of cells.

For a computationally efficient solution to estimate $P\!\left( \boldsymbol{\lambda} \mid \mathbf{X}, \boldsymbol{\theta} \right)$, we use variational inference, an optimization-based approximate Bayesian inference approach.

In brief, we initialize an approximate posterior distribution $q\!\left( \boldsymbol{\lambda} \mid \mathbf{X}, \boldsymbol{\theta} \right)$ that follows a transformed Beta distribution with differentiable parameters describing its shape. $q\!\left( \boldsymbol{\lambda} \mid \mathbf{X}, \boldsymbol{\theta} \right)$ is then optimized to minimize its KL divergence with $P\!\left( \boldsymbol{\lambda} \mid \mathbf{X}, \boldsymbol{\theta} \right)$ by minimizing the evidence lower bound (ELBO) objective function. Additional details of the initialization, generative process, and optimization procedure of $q\!\left( \boldsymbol{\lambda} \mid \mathbf{X}, \boldsymbol{\theta} \right)$ can be viewed in \highlight{Supplementary Methods 2}.

\subsubsection*{Conversion of Bayesian ANOVA parameter posteriors to circadian cosinor parameter posteriors}

Given our approximate posterior $q\!\left( \boldsymbol{\lambda} \mid \mathbf{X}, \boldsymbol{\theta} \right)$, for gene $j$ we can generate a sampled waveform $\lambda_{j}^{*} = \left( \lambda_{j1}^{*},\ \lambda_{j2}^{*}, \lambda_{j3}^{*}, \lambda_{j4}^{*} \right)$, where $\lambda_{j1}^{*},\ \lambda_{j2}^{*},\ \lambda_{j3}^{*},\ \text{and}\ \lambda_{j4}^{*}$ correspond to sampled gene proportions from ZT0, ZT6, ZT12, and ZT18, respectively.

Using this sampled waveform, we apply a Fast Fourier Transform, which yields sampled cosinor parameters corresponding to 12-hour and 24-hour components. We denote the sampled 24-hour circadian components for gene $j$ as $\mu_{j}^{*}$, $A_{j}^{*}$, and $\phi_{j}^{*}$, corresponding to the mesor, amplitude, and acrophase, respectively. Sampling from $q\!\left( \boldsymbol{\lambda} \mid \mathbf{X}, \boldsymbol{\theta} \right)$ repeatedly, we can generate samples of $\mu_{j}, A_{j},$ and $\phi_{j}$ to approximate their posterior distributions.

\subsubsection*{Identification of confident circadian cycling genes}

Confident circadian cyclers were identified using three steps. First, we computed the Bayes factor of each gene using $q\!\left( \boldsymbol{\lambda} \mid \mathbf{X}, \boldsymbol{\theta} \right)$, where 5 $\lambda_{j}^{*}$ samples were drawn for each gene at each time point. Second, we computed the Bayesian evidence of each gene using an altered version of $q\!\left( \boldsymbol{\lambda} \mid \mathbf{X}, \boldsymbol{\theta} \right)$, where the 24-hour circadian component was subtracted out. To do this, 5 ${\lambda_{j}^{*}}^{(\text{flat})}$ samples were drawn for each gene. Samples were then converted to the frequency domain using the Fast Fourier Transform, and the 24-hour circadian component amplitudes were set to zero. Frequency domain samples were then transformed back to the original domain using the Inverse Fast Fourier Transform, and then the Bayesian evidence was computed. Third, we compute a Bayes factor for each gene equal to the Bayesian evidence associated with $q\!\left( \boldsymbol{\lambda} \mid \mathbf{X}, \boldsymbol{\theta} \right)$ divided by the Bayesian evidence associated with the 24-hour component subtracted posterior. We refer to this Bayes factor as the circadian Bayes factor. Confident cycling genes were defined as those with a circadian Bayes factor greater than or equal to 100, unless otherwise stated.

%-----------------------------------------------------------------------
% SUPPLEMENTARY METHODS 1
%-----------------------------------------------------------------------
\section*{Supplementary Methods 1: Estimating the transcript proportion -- dispersion relationship}

We model UMI counts as coming from a Negative Binomial (NB) data likelihood model, as it can account for the overdispersion observed in real count data. While many existing applications of NB to scRNA-seq data model dispersions in a gene-specific fashion, these estimates are known to exhibit high amounts of uncertainty. Methods such as scTransform~\cite{ref134} and DESeq2~\cite{ref135} combat this by an empirical Bayes approach to shrink dispersion estimates, assuming genes with similar means share similar dispersions. We take this a step further and make the simplifying assumption that the transcript proportion--dispersion relationship across all genes is strictly described by a polynomial function $g_{\boldsymbol{\zeta}}\!\left( \lambda_{ij} \right)$ parameterized by $\boldsymbol{\zeta}$. The coefficients of this polynomial, $\boldsymbol{\zeta}$, are estimated as follows:

\subsubsection*{Step 1: Estimate gene-specific proportions and dispersions}

We presume gene counts are distributed according to:

\begin{equation}
X_{ij} \sim \mathrm{NB}\!\left( \widetilde{\lambda_{j}} \cdot L_{i},\, \widetilde{\delta_{j}} \right) \tag{24}
\end{equation}

Where:

\begin{equation}
\mathbb{E}\!\left[ X_{ij} \right] = \widetilde{\lambda_{j}} \cdot L_{i} \tag{25}
\end{equation}

\begin{equation}
\mathrm{Var}\!\left( X_{ij} \right) = \mathbb{E}\!\left[ X_{ij} \right] + \widetilde{\delta_{j}} \left( \mathbb{E}\!\left[ X_{ij} \right] \right)^{2} \tag{26}
\end{equation}

And $\widetilde{\lambda_{j}}$ and $\widetilde{\delta_{j}}$ are the temporary (i.e.\ only used to estimate $\boldsymbol{\zeta}$) gene-specific transcript proportion and dispersion for gene $j$, respectively. Under this likelihood model, for each gene we estimate the maximum likelihood transcript proportion $\widetilde{\lambda_{j}}$ and dispersion $\widetilde{\delta_{j}}$.

\subsubsection*{Step 2: Fit polynomial model describing the transcript proportion -- dispersion relationship across all genes}

Assume $g$ is a polynomial function whose coefficients $\boldsymbol{\zeta}$ defines the transcript $\log$ proportion -- $\log$ dispersion relationship shared across all genes:

\begin{equation}
g_{\boldsymbol{\zeta}}\!\left( \lambda_{ij} \right) = \log\!\left( \delta_{j} \right) = \sum_{k=1}^{|\boldsymbol{\zeta}|} \zeta_{k} \log^{k}\!\left( \lambda_{ij} \right) \tag{27}
\end{equation}

Using OLS, we estimate the coefficients $\boldsymbol{\zeta}$ under the following objective:

\begin{equation}
\underset{\boldsymbol{\zeta}}{\arg\min}\ \sum_{j=1}^{p} \left( \log\!\left( \widetilde{\delta_{j}} \right) - g_{\boldsymbol{\zeta}}\!\left( \widetilde{\lambda_{j}} \right) \right)^{2} \tag{28}
\end{equation}

where $p$ is the number of genes, and $\widetilde{\lambda_{j}}$ and $\widetilde{\delta_{j}}$ are the temporary gene-specific transcript proportion and dispersion for gene $j$ estimated in Step 1.

\subsubsection*{Note}

We note that dispersion estimates may be upwardly biased as a subset of genes cycling over the circadian cycle will be fit under a mean model that assumes flat expression over the circadian cycle. However, we assume the fraction of cycling genes detected is small, and thus should not upwardly bias the dispersion estimates much.

%-----------------------------------------------------------------------
% REFERENCES
%-----------------------------------------------------------------------
\section*{References}

\begin{enumerate}
\item Panza, J. A., Epstein, S. E. \& Quyyumi, A. A. Circadian Variation in Vascular Tone and Its Relation to $\alpha$-Sympathetic Vasoconstrictor Activity. \textit{New England Journal of Medicine} \textbf{325}, 986--990 (1991).

\item Bridges, A. B., McLaren, M., Saniabadi, A., Fisher, T. C. \& Belch, J. J. F. Circadian variation of endothelial cell function, red blood cell deformability and dehydro-thromboxane B2 in healthy volunteers. \textit{Blood Coagulation \& Fibrinolysis} \textbf{2}, 447--452 (1991).

\item Tunçtan, B. \textit{et al.} Circadian variation of nitric oxide synthase activity in mouse tissue. \textit{Chronobiol Int} \textbf{19}, 393--404 (2002).

\item Winter, C. \textit{et al.} Chrono-pharmacological Targeting of the CCL2-CCR2 Axis Ameliorates Atherosclerosis. \textit{Cell Metab} \textbf{28}, 175--182.e5 (2018).

\item He, W. \textit{et al.} Circadian Expression of Migratory Factors Establishes Lineage-Specific Signatures that Guide the Homing of Leukocyte Subsets to Tissues. \textit{Immunity} \textbf{49}, 1175--1190.e7 (2018).

\item Morris, C. J., Purvis, T. E., Hu, K. \& Scheer, F. A. J. L. Circadian misalignment increases cardiovascular disease risk factors in humans. \textit{Proceedings of the National Academy of Sciences} \textbf{113}, E1402--E1411 (2016).

\item Knutsson, A., Jonsson, B. G., Akerstedt, T. \& Orth-Gomer, K. Increased risk of ischaemic heart disease in shift workers. \textit{The Lancet} \textbf{328}, 89--92 (1986).

\item Karlsson, B., Knutsson, A. \& Lindahl, B. Is there an association between shift work and having a metabolic syndrome? Results from a population based study of 27,485 people. \textit{Occup Environ Med} \textbf{58}, 747--752 (2001).

\item Tüchsen, F., Hannerz, H. \& Burr, H. A 12 year prospective study of circulatory disease among Danish shift workers. \textit{Occup Environ Med} \textbf{63}, 451--455 (2006).

\item Lin, C. \textit{et al.} Clock Gene Bmal1 Disruption in Vascular Smooth Muscle Cells Worsens Carotid Atherosclerotic Lesions. \textit{Arterioscler Thromb Vasc Biol} (2022) doi:10.1161/atvbaha.121.316480.

\item Shen, Y. \textit{et al.} BMAL1 modulates smooth muscle cells phenotypic switch towards fibroblast-like cells and stabilizes atherosclerotic plaques by upregulating YAP1. \textit{Biochim Biophys Acta Mol Basis Dis} \textbf{1868}, (2022).

\item Pan, X., Jiang, X. C. \& Hussain, M. M. Impaired cholesterol metabolism and enhanced atherosclerosis in clock Mutant Mice. \textit{Circulation} \textbf{128}, 1758--1769 (2013).

\item Cheng, B. \textit{et al.} Tissue-intrinsic dysfunction of circadian clock confers transplant arteriosclerosis. \textit{Proc Natl Acad Sci U S A} \textbf{108}, 17147--17152 (2011).

\item DuPont, J. J., Kenney, R. M., Patel, A. R. \& Jaffe, I. Z. Sex differences in mechanisms of arterial stiffness. \textit{British Journal of Pharmacology} \textbf{176}, 4208--4225 (2019).

\item Wang, M. \textit{et al.} Gender heterogeneity in dyslipidemia prevalence, trends with age and associated factors in middle age rural Chinese. \textit{Lipids Health Dis} \textbf{19}, (2020).

\item Palmisano, B. T., Zhu, L., Eckel, R. H. \& Stafford, J. M. Sex differences in lipid and lipoprotein metabolism. \textit{Molecular Metabolism} \textbf{15}, 45--55 (2018).

\item Gillis, E. E. \& Sullivan, J. C. Sex Differences in Hypertension: Recent Advances. \textit{Hypertension} \textbf{68}, 1322--1327 (2016).

\item Klein, S. L. \& Flanagan, K. L. Sex differences in immune responses. \textit{Nature Reviews Immunology} \textbf{16}, 626--638 (2016).

\item Man, J. J., Beckman, J. A. \& Jaffe, I. Z. Sex as a Biological Variable in Atherosclerosis. \textit{Circ Res} 1297--1319 (2020) doi:10.1161/CIRCRESAHA.120.315930.

\item Vrijenhoek, J. E. P. \textit{et al.} Sex is associated with the presence of atherosclerotic plaque hemorrhage and modifies the relation between plaque hemorrhage and cardiovascular outcome. \textit{Stroke} \textbf{44}, 3318--3323 (2013).

\item Anderson, S. T. \& FitzGerald, G. A. Sexual dimorphism in body clocks. \textit{Science} \textbf{369}, 1164--1165 (2020).

\item Man, A. W. C., Li, H. \& Xia, N. Circadian rhythm: Potential therapeutic target for atherosclerosis and thrombosis. \textit{International Journal of Molecular Sciences} \textbf{22}, 1--19 (2021).

\item McAlpine, C. S. \& Swirski, F. K. Circadian influence on metabolism and inflammation in atherosclerosis. \textit{Circ Res} \textbf{119}, 131--141 (2016).

\item Paschos, G. K. \& FitzGerald, G. A. Circadian clocks and vascular function. \textit{Circ Res} \textbf{106}, 833--841 (2010).

\item Wolock, S. L., Lopez, R. \& Klein, A. M. Scrublet: Computational Identification of Cell Doublets in Single-Cell Transcriptomic Data. \textit{Cell Syst} \textbf{8}, 281--291.e9 (2019).

\item Lopez, R., Regier, J., Cole, M. B., Jordan, M. I. \& Yosef, N. Deep generative modeling for single-cell transcriptomics. \textit{Nat Methods} \textbf{15}, 1053--1058 (2018).

\item Traag, V. A., Waltman, L. \& van Eck, N. J. From Louvain to Leiden: guaranteeing well-connected communities. \textit{Sci Rep} \textbf{9}, 1--12 (2019).

\item Bennett, M. R., Sinha, S. \& Owens, G. K. Vascular Smooth Muscle Cells in Atherosclerosis. \textit{Circ Res} \textbf{118}, 692--702 (2016).

\item Trigueros-Motos, L. \textit{et al.} Embryological-origin-dependent differences in homeobox expression in adult aorta: Role in regional phenotypic variability and regulation of NF-kB activity. \textit{Arterioscler Thromb Vasc Biol} \textbf{33}, 1248--1256 (2013).

\item Dobnikar, L. \textit{et al.} Disease-relevant transcriptional signatures identified in individual smooth muscle cells from healthy mouse vessels. \textit{Nat Commun} \textbf{9}, (2018).

\item Gao, Y. K. \textit{et al.} A regulator of G protein signaling 5 marked subpopulation of vascular smooth muscle cells is lost during vascular disease. \textit{PLoS One} \textbf{17}, (2022).

\item Thysen, S., Cailotto, F. \& Lories, R. Osteogenesis induced by frizzled-related protein (FRZB) is linked to the netrin-like domain. \textit{Laboratory Investigation} \textbf{96}, 570--580 (2016).

\item Komori, T. Regulation of proliferation, differentiation and functions of osteoblasts by runx2. \textit{International Journal of Molecular Sciences} \textbf{20}, (2019).

\item Lefebvre, V. \& Dvir-Ginzberg, M. SOX9 and the many facets of its regulation in the chondrocyte lineage. \textit{Connective Tissue Research} \textbf{58}, 2--14 (2017).

\item Pan, H. \textit{et al.} Single-Cell Genomics Reveals a Novel Cell State During Smooth Muscle Cell Phenotypic Switching and Potential Therapeutic Targets for Atherosclerosis in Mouse and Human. \textit{Circulation} 2060--2075 (2020) doi:10.1161/circulationaha.120.048378.

\item Zhang, R., Lahens, N. F., Ballance, H. I., Hughes, M. E. \& Hogenesch, J. B. A circadian gene expression atlas in mammals: Implications for biology and medicine. \textit{Proceedings of the National Academy of Sciences} \textbf{111}, 16219--16224 (2014).

\item Wang, Y. \textit{et al.} Clonally expanding smooth muscle cells promote atherosclerosis by escaping efferocytosis and activating the complement cascade. \textit{Proceedings of the National Academy of Sciences} \textbf{117}, 202006348 (2020).

\item Wirka, R. C. \textit{et al.} Atheroprotective roles of smooth muscle cell phenotypic modulation and the TCF21 disease gene as revealed by single-cell analysis. \textit{Nat Med} (2019) doi:10.1038/s41591-019-0512-5.

\item Yap, C., Mieremet, A., de Vries, C. J. M., Micha, D. \& de Waard, V. Six shades of vascular smooth muscle cells illuminated by klf4 (Krüppel-like factor 4). \textit{Arteriosclerosis, Thrombosis, and Vascular Biology} 2693--2707 (2021).

\item Chen, E. Y. \textit{et al.} Enrichr: interactive and collaborative HTML5 gene list enrichment analysis tool. \url{http://amp.pharm.mssm.edu/Enrichr} (2013).

\item Rudic, R. D. \textit{et al.} Bioinformatic analysis of circadian gene oscillation in mouse aorta. \textit{Circulation} \textbf{112}, 2716--2724 (2005).

\item Millar-Craig, M., Bishop, C. \& Raftery, E. B. Circadian variation of blood-pressure. \textit{The Lancet} \textbf{311}, 795--797 (1978).

\item Sur, S. H., Mistlberger, R. E., Morris, M. \& Morris, M. Circadian blood pressure and heart rate rhythms in mice. (1999).

\item Scheer, F. A. J. L. \textit{et al.} The human endogenous circadian system causes greatest platelet activation during the biological morning independent of behaviors. \textit{PLoS One} \textbf{6}, (2011).

\item Luciano, A. K., Santana, J. M., Velazquez, H. \& Sessa, W. C. Akt1 Controls the Timing and Amplitude of Vascular Circadian Gene Expression. \textit{J Biol Rhythms} \textbf{32}, 212--221 (2017).

\item Adamovich, Y., Ladeuix, B., Golik, M., Koeners, M. P. \& Asher, G. Rhythmic Oxygen Levels Reset Circadian Clocks through HIF1$\alpha$. \textit{Cell Metab} \textbf{25}, 93--101 (2017).

\item Manella, G. \textit{et al.} Hypoxia induces a time- and tissue-specific response that elicits intertissue circadian clock misalignment. (2019) doi:10.1073/pnas.1914112117.

\item Paroo, Z., Dipchand, E. S. \& Noble, E. G. Estrogen attenuates postexercise HSP70 expression in skeletal muscle. \textit{Am J Physiol Cell Physiol} \textbf{282}, 245--251 (2002).

\item Beere, H. M. \textit{et al.} Heat-shock protein 70 inhibits apoptosis by preventing recruitment of procaspase-9 to the Apaf-1 apoptosome. \textit{Nature Cell Biology} \textbf{2}, (2000).

\item Lu, T. S. \textit{et al.} Induction of intracellular heat-shock protein 72 prevents the development of vascular smooth muscle cell calcification. \textit{Cardiovasc Res} \textbf{96}, 524--532 (2012).

\item Kim, I. K., Shin, H. M. \& Baek, W. Heat-shock response is associated with decreased production of interleukin-6 in murine aortic vascular smooth muscle cells. \textit{Naunyn Schmiedebergs Arch Pharmacol} \textbf{371}, 27--33 (2005).

\item McCullagh, K. J. A., Cooney, R. \& O'Brien, T. Endothelial nitric oxide synthase induces heat shock protein HSPA6 (HSP70B') in human arterial smooth muscle cells. \textit{Nitric Oxide} \textbf{52}, 41--48 (2016).

\item Shi, J., Yang, Y., Cheng, A., Xu, G. \& He, F. Metabolism of vascular smooth muscle cells in vascular diseases. \textit{Am J Physiol Heart Circ Physiol} \textbf{319}, 613--631 (2020).

\item Zheng, Y., Im, C.-N. \& Seo, J.-S. Inhibitory effect of Hsp70 on angiotensin II-induced vascular smooth muscle cell hypertrophy. \textit{Exp Mol Med} \textbf{38}, 509--518 (2006).

\item Guo, X. Transforming growth factor-$\beta$ and smooth muscle differentiation. \textit{World J Biol Chem} \textbf{3}, 41 (2012).

\item Ku, H. C. \& Cheng, C. F. Master Regulator Activating Transcription Factor 3 (ATF3) in Metabolic Homeostasis and Cancer. \textit{Frontiers in Endocrinology} \textbf{11}, (2020).

\item Yin, H.-M. \textit{et al.} Activating transcription factor 3 coordinates differentiation of cardiac and hematopoietic progenitors by regulating glucose metabolism. \textit{Sci Adv} \textbf{6}, (2020).

\item Miller, C. L. \textit{et al.} Integrative functional genomics identifies regulatory mechanisms at coronary artery disease loci. \textit{Nat Commun} \textbf{7}, (2016).

\item Wang, Y. \textit{et al.} Dynamic changes in chromatin accessibility are associated with the atherogenic transitioning of vascular smooth muscle cells. \textit{Cardiovasc Res} (2021) doi:10.1093/cvr/cvab347.

\item Herring, J. A., Elison, W. S. \& Tessem, J. S. Function of nr4a orphan nuclear receptors in proliferation, apoptosis and fuel utilization across tissues. \textit{Cells} \textbf{8}, (2019).

\item Khambata, R. S., Panayiotou, C. M. \& Hobbs, A. J. Natriuretic peptide receptor-3 underpins the disparate regulation of endothelial and vascular smooth muscle cell proliferation by C-type natriuretic peptide. \textit{Br J Pharmacol} \textbf{164}, 584--597 (2011).

\item Orriols, M. \textit{et al.} Down-regulation of Fibulin-5 is associated with aortic dilation: role of inflammation and epigenetics. \textit{Cardiovasc Res} \textbf{110}, 431--442 (2016).

\item Spencer, J. A. \textit{et al.} Altered vascular remodeling in fibulin-5-deficient mice reveals a role of fibulin-5 in smooth muscle cell proliferation and migration. \textit{Proceedings of the National Academy of Sciences} \textbf{102}, 2946--2951 (2005).

\item Liu, Z. P., Wang, Z., Yanagisawa, H. \& Olson, E. N. Phenotypic modulation of smooth muscle cells through interaction of Foxo4 and Myocardin. \textit{Dev Cell} \textbf{9}, 261--270 (2005).

\item Hartman, R. J. G. \textit{et al.} Sex-Stratified Gene Regulatory Networks Reveal Female Key-Driver Genes of Atherosclerosis Involved in Smooth Muscle Cell Phenotype Switching. \textit{Circulation} 713--726 (2021) doi:10.1161/circulationaha.120.051231.

\item Imamura, T. \textit{et al.} Smad6 inhibits signalling by the TGF-$\beta$ superfamily. \textit{Nature} \textbf{389}, 622--626 (1997).

\item Miyazawa, K. \& Miyazono, K. Regulation of TGF-$\beta$ family signaling by inhibitory smads. \textit{Cold Spring Harb Perspect Biol} \textbf{9}, (2017).

\item Nomiyama, T. \textit{et al.} The NR4A Orphan Nuclear Receptor NOR1 Is Induced by Platelet-derived Growth Factor and Mediates Vascular Smooth Muscle Cell Proliferation. \textit{J Biol Chem} \textbf{281}, (2006).

\item Janich, P. \textit{et al.} The circadian molecular clock creates epidermal stem cell heterogeneity. \textit{Nature} \textbf{480}, 209--214 (2011).

\item Brown, S. A. Circadian clock-mediated control of stem cell division and differentiation: Beyond night and day. \textit{Development (Cambridge)} \textbf{141}, 3105--3111 (2014).

\item Zhang, Z.-B., Sinha, J., Bahrami-Nejad, Z. \& Teruel, M. N. The circadian clock mediates daily bursts of cell differentiation by periodically restricting cell-differentiation commitment. \textit{Proceedings of the National Academy of Sciences} \textbf{119}, (2022).

\item Yeung, C. Y. C. \textit{et al.} Gremlin-2 is a BMP antagonist that is regulated by the circadian clock. \textit{Sci Rep} \textbf{4}, (2014).

\item Zhang, X. \textit{et al.} Sp1 Plays an Important Role in Vascular Calcification Both In Vivo and In Vitro. doi:10.1161/JAHA.117.

\item Lyu, Q. \textit{et al.} Hsp70 and NF-$\kappa$B mediated control of innate inflammatory responses in a canine macrophage cell line. \textit{Int J Mol Sci} \textbf{21}, 1--15 (2020).

\item Cicha, I. \textit{et al.} Connective tissue growth factor is overexpressed in complicated atherosclerotic plaques and induces mononuclear cell chemotaxis in vitro. \textit{Arterioscler Thromb Vasc Biol} \textbf{25}, 1008--1013 (2005).

\item Trizzino, M. \textit{et al.} EGR1 is a gatekeeper of inflammatory enhancers in human macrophages. \textit{Sci Adv} \textbf{7}, (2021).

\item Herring, B. P., Hoggatt, A. M., Griffith, S. L., McClintick, J. N. \& Gallagher, P. J. Inflammation and vascular smooth muscle cell dedifferentiation following carotid artery ligation. \textit{Physiol Genomics} \textbf{49}, 115--126 (2017).

\item Balint, B. \textit{et al.} Collectivization of Vascular Smooth Muscle Cells via TGF-$\beta$--Cadherin-11--Dependent Adhesive Switching. \textit{Arterioscler Thromb Vasc Biol} \textbf{35}, 1254--1264 (2015).

\item Ribeiro-Silva, J. C., Miyakawa, A. A. \& Krieger, J. E. Focal adhesion signaling: Vascular smooth muscle cell contractility beyond calcium mechanisms. \textit{Clinical Science} \textbf{135}, 1189--1207 (2021).

\item He, X. \textit{et al.} Activation of M3AChR (Type 3 Muscarinic Acetylcholine Receptor) and Nrf2 (Nuclear Factor Erythroid 2-Related Factor 2) Signaling by Choline Alleviates Vascular Smooth Muscle Cell Phenotypic Switching and Vascular Remodeling. \textit{Arterioscler Thromb Vasc Biol} \textbf{40}, 2649--2664 (2020).

\item Johnson, P. F. Molecular stop signs: Regulation of cell-cycle arrest by C/EBP transcription factors. \textit{J Cell Sci} \textbf{118}, 2545--2555 (2005).

\item Hauge, S., Macurek, L. \& Syljuåsen, R. G. p21 limits S phase DNA damage caused by the Wee1 inhibitor MK1775. \textit{Cell Cycle} \textbf{18}, 834--847 (2019).

\item Kavurma, M. M. \& Khachigian, L. M. Sp1 inhibits proliferation and induces apoptosis in vascular smooth muscle cells by repressing p21WAF1/Cip1 transcription and cyclin D1-Cdk4-p21WAF1/Cip1 complex formation. \textit{Journal of Biological Chemistry} \textbf{278}, 32537--32543 (2003).

\item Jouffe, C. \textit{et al.} Disruption of the circadian clock component BMAL1 elicits an endocrine adaption impacting on insulin sensitivity and liver disease. \textit{Proceedings of the National Academy of Sciences} \textbf{119}, (2022).

\item Yang, G. \textit{et al.} Timing of expression of the core clock gene Bmal1 influences its effects on aging and survival. \textit{Sci Transl Med} \textbf{8}, 324ra16--324ra16 (2016).

\item McGraw, A. P. \textit{et al.} Aldosterone increases early atherosclerosis and promotes plaque inflammation through a placental growth factor-dependent mechanism. \textit{J Am Heart Assoc} \textbf{2}, (2013).

\item Kiss, M. G. \& Binder, C. J. The multifaceted impact of complement on atherosclerosis. \textit{Atherosclerosis} \textbf{351}, 29--40 (2022).

\item Verdeguer, F. \textit{et al.} Complement regulation in murine and human hypercholesterolemia and role in the control of macrophage and smooth muscle cell proliferation. \textit{Cardiovasc Res} \textbf{76}, 340--350 (2007).

\item Lam, M. T. Y. \textit{et al.} Rev-Erbs repress macrophage gene expression by inhibiting enhancer-directed transcription. \textit{Nature} \textbf{498}, 511--515 (2013).

\item Apostolakis, S. \& Spandidos, D. Chemokines and atherosclerosis: Focus on the CX3CL1/CX3CR1 pathway. \textit{Acta Pharmacologica Sinica} \textbf{34}, 1251--1256 (2013).

\item Poupel, L. \textit{et al.} Pharmacological inhibition of the chemokine receptor, CX3CR1, reduces atherosclerosis in mice. \textit{Arterioscler Thromb Vasc Biol} \textbf{33}, 2297--2305 (2013).

\item Timmons, G. A., O'Siorain, J. R., Kennedy, O. D., Curtis, A. M. \& Early, J. O. Innate Rhythms: Clocks at the Center of Monocyte and Macrophage Function. \textit{Frontiers in Immunology} \textbf{11}, 1743 (2020).

\item Yasumoto, H. \textit{et al.} Dominant negative c-Jun gene transfer inhibits vascular smooth muscle cell proliferation and neointimal hyperplasia in rats. \textit{Gene Therapy} \textbf{8}, (2001).

\item Khachigian, L. M., Fahmy, R. G., Zhang, G., Bobryshev, Y. V. \& Kaniaros, A. c-Jun regulates vascular smooth muscle cell growth and neointima formation after arterial injury. \textit{Journal of Biological Chemistry} \textbf{277}, 22985--22991 (2002).

\item Das, S. \textit{et al.} Regulation of angiotensin II actions by enhancers and super-enhancers in vascular smooth muscle cells. \textit{Nat Commun} \textbf{8}, (2017).

\item Thompson, M. R., Xu, D. \& Williams, B. R. G. ATF3 transcription factor and its emerging roles in immunity and cancer. \textit{Journal of Molecular Medicine} \textbf{87}, 1053--1060 (2009).

\item Herbert, J. M., Lamarche, I., Prabonnaud, V., Dol, F. \& Gauthier, T. Tissue-type plasminogen activator is a potent mitogen for human aortic smooth muscle cells. \textit{Journal of Biological Chemistry} \textbf{269}, 3076--3080 (1994).

\item Ishigami, M., Swertfeger, D. K., Hui, M. S., Granholm, N. A. \& Hui, D. Y. Apolipoprotein E Inhibition of Vascular Smooth Muscle Cell Proliferation but Not the Inhibition of Migration Is Mediated Through Activation of Inducible Nitric Oxide Synthase. \textit{Arterioscler Thromb Vasc Biol} (2000).

\item Lee, J. H. \textit{et al.} Mig-6 gene knockout induces neointimal hyperplasia in the vascular smooth muscle cell. \textit{Dis Markers} \textbf{2014}, (2014).

\item Demirel, E. \textit{et al.} Rgs5 attenuates baseline activity of erk1/2 and promotes growth arrest of vascular smooth muscle cells. \textit{Cells} \textbf{10}, (2021).

\item Daniel, J.-M. \textit{et al.} Regulator of G-Protein Signaling 5 Prevents Smooth Muscle Cell Proliferation and Attenuates Neointima Formation. \textit{Arterioscler Thromb Vasc Biol} \textbf{36}, 317--327 (2016).

\item Ji, Y. \textit{et al.} Pharmacological Targeting of Plasminogen Activator Inhibitor-1 Decreases Vascular Smooth Muscle Cell Migration and Neointima Formation. \textit{Arterioscler Thromb Vasc Biol} \textbf{36}, 2167--2175 (2016).

\item Tong, Y. \textit{et al.} Exosome-mediated transfer of ACE (angiotensin-converting enzyme) from adventitial fibroblasts of spontaneously hypertensive rats promotes vascular smooth muscle cell migration. \textit{Hypertension} \textbf{72}, 881--888 (2018).

\item Huang, C., Santofimia-Castaño, P. \& Iovanna, J. NUPR1: A critical regulator of the antioxidant system. \textit{Cancers} \textbf{13}, (2021).

\item Serrano, R. L., Yu, W., Graham, R. M., Liu-Bryan, R. \& Terkeltaub, R. A vascular smooth muscle cell X-box binding protein 1 and transglutaminase 2 regulatory circuit limits neointimal hyperplasia. \textit{PLoS One} \textbf{14}, (2019).

\item Takaguri, A., Kubo, T., Mori, M. \& Satoh, K. The protective role of YAP1 on ER stress-induced cell death in vascular smooth muscle cells. \textit{Eur J Pharmacol} \textbf{815}, 470--477 (2017).

\item Berlanga, J. J., Baass, A. \& Sonenberg, N. Regulation of poly(A) binding protein function in translation: Characterization of the Paip2 homolog, Paip2B. \textit{RNA} \textbf{12}, 1556--1568 (2006).

\item Lecker, S. H., Goldberg, A. L. \& Mitch, W. E. Protein degradation by the ubiquitin-proteasome pathway in normal and disease states. \textit{Journal of the American Society of Nephrology} \textbf{17}, 1807--1819 (2006).

\item Kobayashi, M. \textit{et al.} The ubiquitin hybrid gene UBA52 regulates ubiquitination of ribosome and sustains embryonic development. \textit{Sci Rep} \textbf{6}, (2016).

\item Zeng, M. \textit{et al.} CRL4Wdr70 regulates H2B monoubiquitination and facilitates Exo1-dependent resection. \textit{Nat Commun} \textbf{7}, (2016).

\item Zeng, M., Tang, Z., Guo, L., Wang, X. \& Liu, C. Wdr70 regulates histone modification and genomic maintenance in fission yeast. \textit{Biochim Biophys Acta Mol Cell Res} \textbf{1867}, (2020).

\item Baeten, J. T. \& Lilly, B. Differential regulation of NOTCH2 and NOTCH3 contribute to their unique functions in vascular smooth muscle cells. \textit{Journal of Biological Chemistry} \textbf{290}, 16226--16237 (2015).

\item Boucher, J. M., Harrington, A., Rostama, B., Lindner, V. \& Liaw, L. A receptor-specific function for Notch2 in mediating vascular smooth muscle cell growth arrest through cyclin-dependent kinase inhibitor 1B. \textit{Circ Res} \textbf{113}, 975--985 (2013).

\item Karakaya, C. \textit{et al.} Notch signaling regulates strain-mediated phenotypic switching of vascular smooth muscle cells. \textit{Front Cell Dev Biol} \textbf{10}, (2022).

\item Schultz, K., Murthy, V., Tatro, J. B. \& Beasley, D. Prolyl hydroxylase 2 deficiency limits proliferation of vascular smooth muscle cells by hypoxia-inducible factor-1-dependent mechanisms. \textit{Am J Physiol Lung Cell Mol Physiol} \textbf{296}, 921--927 (2009).

\item Zhao, Y. \textit{et al.} Rho-associated protein kinase isoforms stimulate proliferation of vascular smooth muscle cells through ERK and induction of cyclin D1 and PCNA. \textit{Biochem Biophys Res Commun} \textbf{432}, 488--493 (2013).

\item Matsumae, H. \textit{et al.} CCN1 knockdown suppresses neointimal hyperplasia in a rat artery balloon injury model. \textit{Arterioscler Thromb Vasc Biol} \textbf{28}, 1077--1083 (2008).

\item Forrest, S. \& McNamara, C. Id family of transcription factors and vascular lesion formation. \textit{Arteriosclerosis, Thrombosis, and Vascular Biology} \textbf{24}, 2014--2020 (2004).

\item Forrest, S. T., Taylor, A. M., Sarembock, I. J., Perlegas, D. \& McNamara, C. A. Phosphorylation regulates Id3 function in vascular smooth muscle cells. \textit{Circ Res} \textbf{95}, 557--559 (2004).

\item Lowery, J. W. \textit{et al.} ID family protein expression and regulation in hypoxic pulmonary hypertension. \textit{Am J Physiol Regul Integr Comp Physiol} \textbf{299}, 1463--1477 (2010).

\item Tejera-Muñoz, A. \textit{et al.} Ccn2 increases tgf-$\beta$ receptor type ii expression in vascular smooth muscle cells: Essential role of ccn2 in the tgf-$\beta$ pathway regulation. \textit{Int J Mol Sci} \textbf{23}, (2022).

\item Chen, C. C. \& Lau, L. F. Functions and mechanisms of action of CCN matricellular proteins. \textit{International Journal of Biochemistry and Cell Biology} \textbf{41}, 771--783 (2009).

\item Suehiro, J. I., Hamakubo, T., Kodama, T., Aird, W. C. \& Minami, T. Vascular endothelial growth factor activation of endothelial cells is mediated by early growth response-3. \textit{Blood} \textbf{115}, 2520--2532 (2010).

\item Liu, D., Evans, I., Britton, G. \& Zachary, I. The zinc-finger transcription factor, early growth response 3, mediates VEGF-induced angiogenesis. \textit{Oncogene} \textbf{27}, 2989--2998 (2008).

\item Phng, L. K. \textit{et al.} Nrarp Coordinates Endothelial Notch and Wnt Signaling to Control Vessel Density in Angiogenesis. \textit{Dev Cell} \textbf{16}, 70--82 (2009).

\item Harjes, U., Bridges, E., McIntyre, A., Fielding, B. A. \& Harris, A. L. Fatty acid-binding protein 4, a point of convergence for angiogenic and metabolic signaling pathways in endothelial cells. \textit{Journal of Biological Chemistry} \textbf{289}, 23168--23176 (2014).

\item Harjes, U. \textit{et al.} Antiangiogenic and tumour inhibitory effects of downregulating tumour endothelial FABP4. \textit{Oncogene} \textbf{36}, 912--921 (2017).

\item Kanemaru, H. \textit{et al.} BATF2 activates DUSP2 gene expression and up-regulates NF-$\kappa$B activity via phospho-STAT3 dephosphorylation. \textit{Int Immunol} \textbf{30}, 255--265 (2018).

\item Jeffrey, K. L. \textit{et al.} Positive regulation of immune cell function and inflammatory responses by phosphatase PAC-1. \textit{Nat Immunol} \textbf{7}, 274--283 (2006).

\item Rigo, A. \textit{et al.} Macrophages may promote cancer growth via a GM-CSF/HB-EGF paracrine loop that is enhanced by CXCL12. \textit{Mol Cancer} \textbf{9}, (2010).

\item Edwards, J. P., Zhang, X. \& Mosser, D. M. The Expression of Heparin-Binding Epidermal Growth Factor-Like Growth Factor by Regulatory Macrophages. \textit{The Journal of Immunology} \textbf{182}, 1929--1939 (2009).

\item Higashiyama, S., Abraham, J. A., Miller, J., Fiddes, J. C. \& Klagsbrun, M. A Heparin-Binding Growth Factor Secreted by Macrophage-Like Cells That Is Related to EGF. \textit{Science} \textbf{251}, 936--939 (1991).

\item Xia, C., Braunstein, Z., Toomey, A. C., Zhong, J. \& Rao, X. S100 proteins as an important regulator of macrophage inflammation. \textit{Frontiers in Immunology} \textbf{8}, (2018).

\item Hafemeister, C. \& Satija, R. Normalization and variance stabilization of single-cell RNA-seq data using regularized negative binomial regression. \textit{Genome Biol} \textbf{20}, 1--15 (2019).

\item Love, M. I., Huber, W. \& Anders, S. Moderated estimation of fold change and dispersion for RNA-seq data with DESeq2. \textit{Genome Biol} \textbf{15}, 1--21 (2014).
\end{enumerate}

%-----------------------------------------------------------------------
% ACKNOWLEDGEMENTS
%-----------------------------------------------------------------------
\section*{Acknowledgements}

This work was supported by 2U54TR001878 (G.A.F), R01GM125301 (M.L.), 5TT32HL007953-22 (B.J.A.), and 2T32HG000046-21 (B.J.A.).

%-----------------------------------------------------------------------
% AUTHOR CONTRIBUTIONS
%-----------------------------------------------------------------------
\section*{Author contributions}

Benjamin J. Auerbach (B.J.A.), Soon Yew Tang (S.Y.T), Ronan Lordan (R.L.), Sean T. Anderson (S.T.A.), Mingyao Li (M.L.), and Garret A. FitzGerald (G.A.F.) contributed to the project. The study was conceived of by G.A.F. and B.J.A. S.Y.T., R.L., and S.T.A. performed the mouse experiments and generated the scRNA-seq data from the mouse aorta. B.J.A. performed the data analysis, with input from M.L and G.A.F. B.J.A. wrote the paper with feedback from M.L. and G.A.F.

%-----------------------------------------------------------------------
% COMPETING INTERESTS
%-----------------------------------------------------------------------
\section*{Competing financial interests}

G.A.F is a Senior Advisor to Calico Laboratories. The remaining authors declare no competing financial interests. All authors declare no competing non-financial interests.

%-----------------------------------------------------------------------
% SOFTWARE & DATA AVAILABILITY
%-----------------------------------------------------------------------
\section*{Software availability}

Scripts used to analyze the data and generate the figures can be found at \url{https://github.com/bauerbach95/misalignment}.

\section*{Data availability}

Raw sequencing and UMI count matrices were deposited to the Gene Expression Omnibus under accession code \highlight{GSEXXXX}.

\end{document}